% Aeolion technical report: the DAVE-ML flight-model tabulation.
% Build:  pdflatex report && pdflatex report   (twice, for the ToC and
% cross-references; no bibtex -- the bibliography is a thebibliography
% environment, and latexmk does not run on the build machine, no perl.)
%
% The data tables in tables/ are GENERATED from the solver's own JSON by
% make-tables.py; never hand-edit them (the house rule that keeps every
% published number traceable to a run).
\documentclass[11pt,a4paper]{article}

\usepackage[margin=25mm]{geometry}
% T1 so the Turkish breve/cedilla accents in the author names and the
% bibliography render under pdflatex. Names are written as LaTeX accent
% commands rather than literal UTF-8, matching the papers' own preamble
% note.
\usepackage[T1]{fontenc}
\usepackage{amsmath}
\usepackage{amssymb}
\usepackage{booktabs}
\usepackage{longtable}
\usepackage{array}
\usepackage{graphicx}
\usepackage{xcolor}
\usepackage{listings}
\usepackage[hidelinks]{hyperref}
\usepackage{caption}

\lstset{
  basicstyle=\ttfamily\footnotesize,
  breaklines=true,
  frame=single,
  rulecolor=\color{black!30},
  backgroundcolor=\color{black!3},
  columns=fullflexible,
}

\newcommand{\varid}[1]{\texttt{\small #1}}
\newcommand{\dg}{^\circ}
\newcommand{\nomen}[2]{\textit{#1} & \quad #2 \\}

\title{\textbf{A DAVE-ML Flight-Model Tabulation for the Aetherion
Ducted-Fan Tail-Sitter}\\[2mm]
\large Structure, Conventions, Provenance, and Declared Limits}

\author{
  Onur Tuncer\thanks{Professor, Department of Aeronautical Engineering,
    Istanbul Technical University; onur.tuncer@itu.edu.tr}
  \and
  \c{C}a\u{g}lar \"U\c{c}ler\thanks{Associate Professor,
    \"Ozye\u{g}in University; caglar.ucler@ozyegin.edu.tr}
  \and
  Ahmet G\"une\c{s}\thanks{Associate Professor, Istanbul Technical
    University}
}
\date{15 August 2026 \\[1mm] \small Aeolion Technical Report ATR--2026--01}

\begin{document}
\maketitle

\begin{abstract}
\noindent
This report specifies a single-file DAVE-ML flight model---the encoding
of the American National Standard ANSI/AIAA S-119-2011---for the
Aetherion ducted-fan tail-sitter, assembled from the Aeolion aerodynamic
toolkit. It states the model's structure, its reference-frame
and sign conventions, the nondimensionalisation and buildup of every
tabulated quantity, the provenance of each of its 36 tables, and---given
equal weight---the boundaries of what the underlying methods can support.
Three structural decisions carry the design. First, the model is one file
rather than the conventional aerodynamics/propulsion pair, because the
power-on interaction increments depend on breakpoints owned by both
halves and DAVE-ML provides no cross-file reference mechanism. Second,
forward speed and rotor speed are collapsed into the advance ratio $J$,
which is exact for the inviscid rotating-frame solve that generates the
propulsor tables. Third, the model declares \emph{no} Mach and \emph{no}
Reynolds axis: the toolkit is incompressible throughout and every source
sweep was run at a single flight condition, so such axes would encode a
fidelity that does not exist. The tabulation is verified not by
inspection but by generated \texttt{checkData} static shots evaluated
through an in-repo interpreter, including a moment-transfer shot that
would detect the specific frame-conversion error that has occurred once
in this codebase's history. Two blocks of the model are declared
incomplete and are identified as such in the file itself: the
fan-on-airframe interaction increments await an upstream-induction model,
and the propulsor tables are axial-inflow only.
\end{abstract}

\tableofcontents
\newpage

%======================================================================
\section*{Nomenclature}
\addcontentsline{toc}{section}{Nomenclature}

\begin{longtable}{@{}p{28mm}@{}p{112mm}@{}}
\nomen{A\!R}{= wing aspect ratio, $b^2/S$}
\nomen{b}{= wing reference span, m}
\nomen{$\bar c$}{= wing reference chord, m}
\nomen{$C_D$}{= total drag coefficient}
\nomen{$C_{D_0}$}{= body and duct parasite drag coefficient}
\nomen{$C_{D_i}$}{= induced drag coefficient}
\nomen{$c_d$}{= section profile drag coefficient}
\nomen{$C_L$}{= lift coefficient}
\nomen{$C_{L\max}$}{= maximum lift coefficient}
\nomen{$C_l$, $C_m$, $C_n$}{= rolling, pitching, yawing moment coefficients, body axes}
\nomen{$C_{l_p}$, $C_{l_r}$}{= roll damping and roll-due-to-yaw-rate derivatives, per $\hat p$, $\hat r$}
\nomen{$C_{m_q}$}{= pitch damping derivative, per $\hat q$}
\nomen{$C_{n_p}$, $C_{n_r}$}{= yaw-due-to-roll-rate and yaw damping derivatives}
\nomen{$C_N$}{= normal-force coefficient, chord-plane normal}
\nomen{$C_Q$}{= rotor torque coefficient, $Q/\rho n^2 D^5$}
\nomen{$C_T$}{= rotor thrust coefficient, $T/\rho n^2 D^4$}
\nomen{$C_X$, $C_Y$, $C_Z$}{= body-axis force coefficients}
\nomen{$C_{d\max}$}{= deep-stall section drag maximum (Viterna anchor)}
\nomen{$D$}{= rotor diameter, m}
\nomen{$D_i$}{= induced drag, N}
\nomen{$f$}{= Kirchhoff attached-flow fraction, chordwise separation point}
\nomen{$F_x$, $F_y$, $F_z$}{= total body-axis forces, N}
\nomen{$J$}{= advance ratio, $V_\infty/nD$}
\nomen{$M$}{= Mach number}
\nomen{$M_x$, $M_y$, $M_z$}{= total body-axis moments about the reference point, N\,m}
\nomen{$n$}{= rotor rotational speed, rev/s}
\nomen{$p$, $q$, $r$}{= body-axis roll, pitch, yaw rates, rad/s}
\nomen{$\hat p$, $\hat q$, $\hat r$}{= reduced (nondimensional) body rates}
\nomen{$\bar q$}{= free-stream dynamic pressure, $\tfrac12 \rho V_\infty^2$, Pa}
\nomen{$Q$}{= rotor shaft torque, N\,m}
\nomen{$\mathrm{Re}_n$}{= leading-edge-normal Reynolds number}
\nomen{$r_{\mathrm{LE}}$}{= section leading-edge radius, m}
\nomen{$S$}{= wing gross planform reference area, m$^2$}
\nomen{$s$}{= swirl factor of the rotor--vane momentum budget}
\nomen{$T$}{= rotor thrust, N}
\nomen{$T_c$}{= airframe-normalised thrust coefficient, $T/\bar q S$}
\nomen{$\hat u$}{= free-stream unit vector, body axes}
\nomen{$V_\infty$}{= true airspeed, m/s}
\nomen{$v_i$}{= axial induced velocity, m/s}
\nomen{$x$, $y$, $z$}{= body-axis coordinates, m}
\nomen{$x_{\mathrm{cp}}$}{= centre of pressure, fraction of chord}
\nomen{$\alpha$}{= angle of attack, deg}
\nomen{$\alpha_{\mathrm{disk}}$}{= angle between the free stream and the rotor axis, deg}
\nomen{$\beta$}{= sideslip angle, deg}
\nomen{$\delta_a$}{= aileron deflection (right surface), deg}
\nomen{$\delta_P$, $\delta_Y$, $\delta_R$}{= vane pitch, yaw, roll mode commands, deg}
\nomen{$\delta_{\mathrm{b}}$, $\delta_{\mathrm{l}}$,}{}
\nomen{$\delta_{\mathrm{t}}$, $\delta_{\mathrm{r}}$}{= individual vane deflections (bottom, left, top, right), deg}
\nomen{$\eta$}{= span fraction (wing) or exit-radius fraction (vanes)}
\nomen{$\rho$}{= air density, kg/m$^3$}
\nomen{$\sigma$}{= total inclination of the free stream to the chord plane, $\sin\sigma = \sin\alpha\cos\beta$, deg}
\nomen{$\phi_w$}{= crossflow azimuth about the rotor axis, deg}
\nomen{$\Omega$}{= rotor angular speed, rad/s}
\end{longtable}

\noindent\textit{Subscripts}
\begin{longtable}{@{}p{28mm}@{}p{112mm}@{}}
\nomen{aero}{airframe, power-off contribution}
\nomen{cpl}{fan-on-airframe interaction contribution}
\nomen{disk}{referred to the rotor disk}
\nomen{frd}{contract body frame: $x$ forward, $y$ right, $z$ down}
\nomen{mrp}{moment reference point}
\nomen{prop}{propulsor contribution}
\nomen{vlm}{solver frame: $x$ aft, $y$ right, $z$ up}
\nomen{$\infty$}{free-stream condition}
\end{longtable}

\noindent\textit{Abbreviations}
\begin{longtable}{@{}p{28mm}@{}p{112mm}@{}}
\nomen{BP}{breakpoint set}
\nomen{CST}{class/shape transformation airfoil parameterisation}
\nomen{DAVE-ML}{Dynamic Aerospace Vehicle Exchange Markup Language (AIAA S-119)}
\nomen{FRD}{forward--right--down body axes}
\nomen{MRP}{moment reference point}
\nomen{RMS}{root mean square}
\nomen{VLM}{vortex-lattice method}
\end{longtable}

%======================================================================
\section{Purpose and scope}

The Aeolion toolkit computes the aerodynamics of a ducted-fan tail-sitter
by a family of related potential-flow methods: a vortex-lattice solve
with source-panelled bodies, a sectional viscous coupling, an anchored
post-stall section model, and a rotating-frame lattice for the propulsor
with a two-way rotor--vane coupling. These methods answer questions one
attitude at a time. A flight simulator, a trim routine, or a control
allocator instead needs the whole envelope available at once, evaluated
in microseconds, in a form that no longer depends on the solver being
installed.

DAVE-ML---the Dynamic Aerospace Vehicle Exchange Markup Language, the
encoding defined by the American National Standard ANSI/AIAA
S-119-2011~\cite{s119} and its Annex B reference
document~\cite{davemlref}---is the standard interchange for exactly this:
gridded tables, their breakpoints, the algebraic buildup that assembles
them, and, importantly, machine-checkable verification data carried in
the same document. This report specifies the Aetherion model in that
language.

\subsection{Standards conformance}
\label{sec:standards}

The model is written against the following documents, and departures
from any of them are stated where they occur.

\begin{table}[htbp]
\centering
\caption{Governing standards.}
\label{tab:standards}
\begin{tabular}{@{}llp{62mm}@{}}
\toprule
\textbf{Document} & \textbf{Ref.} & \textbf{Governs} \\
\midrule
ANSI/AIAA S-119-2011 & \cite{s119} & the exchange standard itself;
Annex A standard variable names, Annex B the DAVE-ML reference \\
DAVE-ML Reference 2.0.1 & \cite{davemlref} & element syntax, the
\texttt{DAVEfunc} DTD, \texttt{checkData} semantics \\
ANSI/AIAA R-004-1992 & \cite{r004} & body-axis definitions, angle and
force/moment sign conventions \\
ISO 1151-1:1988 & \cite{iso1151} & flight-dynamics concepts, quantities
and symbols; the basis of R-004 \\
XML 1.1, MathML 2.0 & \cite{xml,mathml} & the two open standards
DAVE-ML is built upon \\
\bottomrule
\end{tabular}
\end{table}

The body frame of Section~\ref{sec:frames}---$x$ forward, $y$ toward the
starboard wing, $z$ down, with moments positive by the right-hand rule
about those axes---is the standard aeronautical body axis system of
ANSI/AIAA R-004-1992~\cite{r004}, itself an adaptation of ISO
1151-1:1988~\cite{iso1151}. The reduced-rate nondimensionalisation of
Eq.~\eqref{eq:reduced} and the moment coefficient reference lengths
follow the same document. Adopting the standard frame is what makes the
solver-side conversion of Eq.~\eqref{eq:flip} necessary rather than
optional: the solver's internal axes are not the standard ones, so the
boundary between them must be explicit.

\paragraph{Conformance of the variable names.} S-119 Annex A prescribes
standard variable names for common flight-dynamics quantities, and this
model's identifiers (Table~\ref{tab:inputs}) were chosen instead for
namespacing---the \varid{aero*}, \varid{prop*}, \varid{coupling*}
grouping described below. The two requirements turn out not to conflict,
because the standard separates the two roles that a single identifier
would otherwise have to serve. A \texttt{variableDef} carries both a
\texttt{varID}, ``an internal identifier that is unique within the
file'' and unconstrained in form, and a \texttt{name}, which ``should
correspond to the standard AIAA parameter name''~\cite{davemlref}. The
model therefore keeps its namespaced \texttt{varID}s for the structural
reason above and carries the Annex A name alongside, so a consumer
expecting standard names finds them without the file losing its
split-ability.

\paragraph{Three attributes the model must populate.} The same element
provides declarations for exactly the ambiguities
Section~\ref{sec:frames} is concerned with, and they are more useful than
prose:

\begin{itemize}
\item \texttt{axisSystem} names the frame a signal is measured in
(body, inertial, \ldots). Every force, moment and rate in this model is
body-axis, and says so in the attribute rather than only in this report.
\item \texttt{sign} declares the positive direction with a short token
---\texttt{+UP}, \texttt{+RWD}, \texttt{TED} for trailing-edge-down.
This refines the position taken in Section~\ref{sec:frames}. The
objection there is to a \emph{prose gloss standing in place of} a
verified convention, not to a machine-readable declaration: a
\texttt{sign} attribute is checkable by tooling and is exactly where the
convention belongs. The generated static shot remains the numerical
proof; the attribute is the statement of intent, and a mismatch between
them is itself a defect worth catching.
\item \texttt{symbol} carries the conventional symbol, letting a
consumer's documentation render $\alpha$ rather than
\varid{alphaDeg}.
\end{itemize}

\noindent
The \texttt{alias} attribute---a facility-specific name outside the AIAA
convention---is available and is deliberately not used, the reference
itself discouraging it for portability.

The report covers the model's structure and conventions, the complete
set of variables, equations and tables, the provenance of each table, and
the declared limits. It does not reproduce the underlying theory, which
is documented in \texttt{doc/theory.rst} and in the three journal
manuscripts the tables draw on~\cite{partone,parttwo,scitech}.

\subsection{What is deliberately absent}

A tabulated model is a claim about what the generating method knows.
Three axes a reader might expect are absent by decision, and their
absence is a statement of honesty rather than an omission of work:

\begin{itemize}
\item \textbf{No Mach axis.} Every solver in the toolkit is
incompressible. Mach number is accepted by the section-model interfaces
and explicitly ignored; there is no Prandtl--Glauert correction anywhere
in the chain. A Mach breakpoint would therefore be a fabricated axis: the
generator would produce identical numbers at every breakpoint, and a
consumer would silently trust the resulting flat interpolation as
physics. Validity $M < 0.3$ is declared in the file header instead.
\item \textbf{No Reynolds axis.} All airframe sweeps were run at a single
flight condition ($V_\infty = 25$~m/s, $\mathrm{Re}_n \approx
3\times10^5$). The separation behaviour that shapes the post-stall tables
was computed at that condition and at no other. The reference condition
is declared; an axis is not invented.
\item \textbf{No disk-incidence axis on the propulsor.} See
Section~\ref{sec:axialonly}. The rotor--vane machinery is axisymmetric
end to end, so non-axial inflow is not representable without new
modelling.
\end{itemize}

%======================================================================
\section{Model architecture}

\subsection{One file, not two}

The conventional split places airframe aerodynamics in one DAVE-ML file
and propulsion in another. This model does not, for a specific reason:
the power-on interaction increments $\Delta C(\alpha, T_c)$ are indexed
by an airframe breakpoint ($\alpha$) and a propulsion-derived state
($T_c$), and DAVE-ML has no mechanism by which one file may reference
another's variables or breakpoint sets. Split across two files, the
interaction tables' input contract would be enforced only by naming
discipline between two documents that version independently.

A single file additionally permits the total-wrench buildup of
Section~\ref{sec:buildup}---the dimensionalisation of each coefficient
group by its own reference quantity and the summation of the three
contributions---to live inside the exchanged model as MathML
\texttt{calculation} elements. This moves the most error-prone step in
any consumer integration from the consumer's code into territory that the
file's own verification data can pin.

The cost of one file is coarser regeneration granularity, and it is
avoided by construction: each sweep driver exports its own JSON, and a
separate assembler builds the XML from the cached JSONs
(Section~\ref{sec:pipeline}). Regenerating the vane map does not rerun
the post-stall map. Variable identifiers are namespaced
\varid{aero*}, \varid{prop*}, and \varid{coupling*}, so that a split into
separate files---should a consumer ever require it---is a mechanical
transformation rather than a redesign.

\subsection{Structure at a glance}

The file carries, in order: a header declaring provenance, reference
condition, and validity envelope; the input variable definitions
(Table~\ref{tab:inputs}); the constant definitions
(Table~\ref{tab:constants}); the derived variables as MathML calculations
(Section~\ref{sec:derived}); the breakpoint sets
(Table~\ref{tab:breakpoints}); the 36 gridded table definitions grouped
by namespace (Section~\ref{sec:inventory}); the buildup
(Section~\ref{sec:buildup}); and the \texttt{checkData} static shots
(Section~\ref{sec:verification}).

%======================================================================
\section{Reference frames, signs, and the moment reference}
\label{sec:frames}

This section is normative. Every quantity in the file obeys it, and the
generator implements it at exactly two code sites.

\subsection{The two frames}

The geometry contract states its own frame once and requires the consumer
to convert explicitly at ingest. That frame,
\texttt{aetherion\_body\_frd}, is $x$ forward, $y$ right, $z$ down---the
standard flight-dynamics body frame, and the frame of every quantity in
the DAVE-ML file. The solver works internally in $x$ aft, $y$ right, $z$
up. The two are related by a rotation of $180\dg$ about $y$:
\begin{equation}
x_{\mathrm{frd}} = -x_{\mathrm{vlm}}, \qquad
y_{\mathrm{frd}} = \phantom{-}y_{\mathrm{vlm}}, \qquad
z_{\mathrm{frd}} = -z_{\mathrm{vlm}}.
\label{eq:flip}
\end{equation}

Equation~\eqref{eq:flip} is a \emph{proper} rotation. Two consequences
follow, and both matter. First, handedness is preserved, so a rotation
axis (a pseudovector) transforms exactly as an ordinary vector does, and
the right-hand rule about a converted axis still means what it meant
before conversion. Second, forces and moments transform identically under
it, giving Table~\ref{tab:signs}.

\begin{table}[htbp]
\centering
\caption{Behaviour of each wrench component under the frame conversion of
Eq.~\eqref{eq:flip}.}
\label{tab:signs}
\begin{tabular}{@{}lll@{}}
\toprule
Component & Transformation & Physical quantity \\
\midrule
$C_X$, $C_Z$ & sign flips & axial and normal force \\
$C_Y$ & invariant & side force \\
$C_l$, $C_n$ & sign flips & rolling and yawing moment \\
$C_m$ & invariant & pitching moment \\
\bottomrule
\end{tabular}
\end{table}

The generator realises the conversion as
\begin{equation}
F_x = -D_i, \quad F_y = Y, \quad F_z = -L, \qquad
M_x = -M_x^{\mathrm{vlm}}, \quad M_y = M_y^{\mathrm{vlm}}, \quad
M_z = -M_z^{\mathrm{vlm}},
\label{eq:export}
\end{equation}
where $D_i$, $Y$, $L$ are the solver's streamwise, lateral and normal
force components.

\subsection{A recorded failure, and the check that now catches it}

The conversion in Eq.~\eqref{eq:flip} applies to the moment reference
\emph{point} as much as to the wrench. Omitting it there does not produce
an obviously wrong answer; it places the reference point an equal
distance on the wrong side of the origin, introducing a spurious moment
arm of twice the offset. This occurred once in this codebase: a sweep
driver passed the contract's moment reference point into the solver
without the flip, producing a moment-arm error of approximately
$0.48$~m---about one body length---and inflating $C_{m\alpha}$ by an
order of magnitude. The error was invisible to every existing test,
because no test compared moments across two reference points.

The model therefore adopts two structural defences. First, the generator
performs the conversion at exactly two sites---reference-point ingest and
wrench export, Eq.~\eqref{eq:export}---and nowhere else, so that there is
one place to audit rather than a scattering of sign flips. Second, the
file carries a \emph{moment-transfer static shot}: the same flight
condition exported about the contract's reference point and about a point
displaced $+0.10$~m in $x_{\mathrm{frd}}$. The difference in $C_m$ must
equal the hand-derivable transfer arm,
\begin{equation}
\Delta C_m = \frac{\Delta x \, F_z - \Delta z \, F_x}{\bar q S \bar c},
\label{eq:transfer}
\end{equation}
evaluated at $\Delta x = 0.10$~m, $\Delta z = 0$. A wrong flip fails this
shot by $2 \times 0.2423$~m of arm, roughly two reference chords, which
cannot be mistaken for numerical noise.

\subsection{Moment reference point}

All moments in the file are taken about the contract's stated moment
reference point,
\begin{equation}
(x, y, z)_{\mathrm{mrp}} = (-0.2423,\ 0,\ 0)~\mathrm{m}, \qquad
\text{contract frame (FRD)},
\end{equation}
declared twice in the document: as prose in the file header, and as three
constant variable definitions \varid{XmrpM}, \varid{YmrpM},
\varid{ZmrpM}. The redundancy is deliberate---a consumer transferring the
model to its own centre of gravity needs the point as data, not as a
comment, and the transfer is then Eq.~\eqref{eq:transfer} generalised to
all three axes:
\begin{align}
M_x' &= M_x - \Delta y\,F_z + \Delta z\,F_y, \\
M_y' &= M_y - \Delta z\,F_x + \Delta x\,F_z, \\
M_z' &= M_z - \Delta x\,F_y + \Delta y\,F_x,
\end{align}
with $\Delta = \mathbf{r}_{\mathrm{new}} - \mathbf{r}_{\mathrm{mrp}}$.

\subsection{Deflection sign convention}

Positive deflection of any control surface is the right-hand rule about
that surface's hinge-axis vector \emph{as stated in the contract's FRD
frame}. This is the entire normative statement. English glosses
(``positive is trailing edge down'') are excluded from the file's
normative text, because such a description must be re-derived in each
frame and a re-derivation that goes wrong inverts half the controls
silently. The numeric meaning of each sign is instead pinned by a
generated static shot, and each control's variable definition cites the
shot that defines it.

%======================================================================
\section{Nondimensionalisation}
\label{sec:nondim}

Three normalisation groups appear in the model, and mixing them is the
most likely consumer error---which is why the buildup of
Section~\ref{sec:buildup} is carried inside the file.

\paragraph{Airframe coefficients} use free-stream dynamic pressure and
wing reference geometry:
\begin{equation}
C_X = \frac{F_x^{\mathrm{aero}}}{\bar q S}, \quad
C_Y = \frac{F_y^{\mathrm{aero}}}{\bar q S}, \quad
C_Z = \frac{F_z^{\mathrm{aero}}}{\bar q S}, \quad
C_l = \frac{M_x^{\mathrm{aero}}}{\bar q S b}, \quad
C_m = \frac{M_y^{\mathrm{aero}}}{\bar q S \bar c}, \quad
C_n = \frac{M_z^{\mathrm{aero}}}{\bar q S b},
\end{equation}
with
\begin{equation}
\bar q = \tfrac{1}{2}\rho V_\infty^2 .
\end{equation}

\paragraph{Propulsor coefficients} use the propeller convention, which
remains finite at zero airspeed and is therefore the only viable choice
for a tail-sitter that must be modelled in hover:
\begin{equation}
C_T = \frac{T}{\rho n^2 D^4}, \qquad
C_Q = \frac{Q}{\rho n^2 D^5},
\label{eq:propnorm}
\end{equation}
with $n$ in revolutions per second. The declared validity excludes
$n \to 0$, where this group degenerates; windmilling and
low-rotor-speed descent are outside the model's envelope and are
identified as future work.

\subsection{Drag accounting}
\label{sec:drag}

Drag deserves its own statement, because the body-axis force tables are
routinely misread as carrying induced drag alone, and because one of the
three physical contributions is not computed anywhere in the toolkit.

\begin{table}[htbp]
\centering
\caption{The three drag contributions, where each lives, and its status.}
\label{tab:drag}
\begin{tabular}{@{}llp{52mm}@{}}
\toprule
\textbf{Contribution} & \textbf{Carried by} & \textbf{Status} \\
\midrule
Induced & \varid{aeroCX}, \varid{aeroCZ} & computed (near-field
Kutta--Joukowski; cross-checked against a Trefftz-plane integral) \\
Wing section profile & the same tables & computed---the section model's
$c_d$ is integrated per strip and folded into the same force vector \\
Body and duct parasite & \varid{aeroCD0} & computed by a component
buildup plus a crossflow branch, Eq.~\eqref{eq:cd0} \\
\bottomrule
\end{tabular}
\end{table}

The first point is that the airframe force tables are \emph{not} induced-drag
tables. The Level-2 coupled solve accumulates the circulatory
(Kutta--Joukowski) force and the sectional profile force into one total
per strip, so the exported body-axis coefficients already contain the
wing's viscous drag as the section model computes it. The attached-flow
study~\cite{partone} reports induced drag alone, both near-field and by a
Trefftz-plane integral; the post-stall study~\cite{parttwo} reports
induced plus profile. Neither reports a complete drag polar.

The second point is the contribution the solves cannot give. A
potential-flow method cannot produce the parasite drag of a closed body:
source-panelled bodies carry zero net force in uniform flow to machine
precision, which is d'Alembert's paradox behaving correctly rather than a
defect. That contribution must come from a component buildup---skin
friction with form and interference factors, in the manner of Raymer and
Hoerner~\cite{raymer2018,hoerner1965fluiddynamicdrag}---driven by wetted
areas assembled from the geometry contract.

Two consequences shape the model, and both are easy to get wrong:

\begin{enumerate}
\item \textbf{A whole-airframe $C_{D_0}$ must never be added on top of
these tables.} The wing's skin friction is already inside the section
$c_d$. The parasite term admitted to this model is \emph{body and duct
wetted area only}; adding a conventional airframe buildup would
double-count the wing.
\item \textbf{Parasite drag cannot be a constant here.} The tables span
$\alpha$ to $90\dg$. A slender body broadside to the flow carries a
crossflow drag of order unity on its projected area---comparable to the
entire tabulated normal force---and that contribution varies strongly
with attitude. A single constant $C_{D_0}$, the conventional choice for a
cruise-envelope model, would underpredict axial force in precisely the
deep-stall regime this model exists to cover.
\end{enumerate}

Parasite drag is therefore tabulated as \varid{aeroCD0}$(\alpha)$: a
component buildup for the friction branch plus a slender-body crossflow
term,
\begin{equation}
C_{D_0}(\alpha) = C_{D_0,\mathrm{friction}}
 + \eta\,C_{d_c}\,\frac{S_{\mathrm{plan}}}{S}\,\lvert\sin^3\alpha\rvert ,
\label{eq:cd0}
\end{equation}
following Allen and Perkins and Jorgensen~\cite{allenPerkins1951,
jorgensen1977}, with $S_{\mathrm{plan}}$ the body's side-projected area,
$C_{d_c} = 1.2$ the subcritical circular-cylinder crossflow drag
coefficient, and $\eta = 0.65$ the finite-length proportionality factor
at this body's fineness ratio of 5. The measured values are in
Section~\ref{sec:parasiteresults}.

In the assembled model the parasite term is resolved into body axes along
the velocity direction and added to the force buildup of
Section~\ref{sec:buildup}, so that the total drag
\begin{equation}
C_D = \underbrace{C_{D_i} + C_{d,\mathrm{wing}}}_{\text{inside }
\varid{aeroCX},\, \varid{aeroCY},\, \varid{aeroCZ}} \;+\;
\underbrace{\varid{aeroCD0}(\alpha)}_{\text{body{+}duct}}
\end{equation}
is recovered from the assembled body-axis components by the projection
onto the free-stream direction of Eq.~\eqref{eq:uhat},
\begin{equation}
C_D = -\bigl(C_X \cos\alpha \cos\beta + C_Y \sin\beta
+ C_Z \sin\alpha\cos\beta\bigr),
\label{eq:CDrecover}
\end{equation}
with no further term: once the buildup has run, every drag contribution
is already inside the body-axis components, and adding \varid{aeroCD0}
to Eq.~\eqref{eq:CDrecover} would count it twice.

\paragraph{Rate derivatives} use the conventional reduced-rate
nondimensionalisation,
\begin{equation}
\hat p = \frac{pb}{2V_\infty}, \qquad
\hat q = \frac{q\bar c}{2V_\infty}, \qquad
\hat r = \frac{rb}{2V_\infty},
\label{eq:reduced}
\end{equation}
so that $C_{l_p} = \partial C_l / \partial \hat p$ and so forth. This
convention is worth stating explicitly because its reciprocal was once
implemented in this codebase: the reduced-rate fields multiplied by
$b/2V_\infty$ where Eq.~\eqref{eq:reduced} requires division by that
factor---an error of $(2V_\infty/b)^2 \approx 2000$ at these
numbers, which reported $C_{l_p} = -2\times 10^{-4}$ for a wing whose
textbook value is $-0.45$. With the convention corrected the solver
returns $-0.452$, and that value is now pinned both by a unit test and by
a static shot in this file.

%======================================================================
\section{Variables}
\label{sec:variables}

\subsection{Inputs}

\begin{longtable}{@{}llll@{}}
\caption{Input variable definitions. The \texttt{varID} is namespaced for
this file's structure; the \texttt{name} is the Annex A standard
parameter name a consumer will look for. Angles are in degrees
throughout, which the \texttt{units} attribute states explicitly.}
\label{tab:inputs} \\
\toprule
\textbf{varID} & \textbf{name} (Annex A) & \textbf{Units} &
\textbf{Symbol} \\
\midrule
\endfirsthead
\multicolumn{4}{@{}l}{\small\itshape Table~\ref{tab:inputs}, continued} \\
\toprule
\textbf{varID} & \textbf{name} (Annex A) & \textbf{Units} &
\textbf{Symbol} \\
\midrule
\endhead
\bottomrule
\endfoot
\varid{alphaDeg} & \varid{angleOfAttack} & deg & $\alpha$ \\
\varid{betaDeg} & \varid{angleOfSideslip} & deg & $\beta$ \\
\varid{trueAirspeedMps} & \varid{trueAirspeed} & m/s & $V_\infty$ \\
\varid{airDensityKgpm3} & \varid{airDensity} & kg/m$^3$ & $\rho$ \\
\varid{rollRateRadps} & \varid{bodyAngularRate\_Roll} & rad/s & $p$ \\
\varid{pitchRateRadps} & \varid{bodyAngularRate\_Pitch} & rad/s & $q$ \\
\varid{yawRateRadps} & \varid{bodyAngularRate\_Yaw} & rad/s & $r$ \\
\varid{propSpeedRevps} & \varid{propellerSpeed} & rev/s & $n$ \\
\varid{aileronDeg} & \varid{aileronDeflection} & deg & $\delta_a$ \\
\varid{vanePitchDeg} & \varid{vaneDeflection\_Pitch} & deg & $\delta_P$ \\
\varid{vaneYawDeg} & \varid{vaneDeflection\_Yaw} & deg & $\delta_Y$ \\
\varid{vaneRollDeg} & \varid{vaneDeflection\_Roll} & deg & $\delta_R$ \\
\end{longtable}

\noindent
Every entry additionally carries \texttt{axisSystem="body"} where a frame
applies, and a \texttt{sign} token: \varid{aileronDeg} is the right
surface's angle with the left slaved to its negative, and the vane
commands resolve to per-vane angles through Eq.~\eqref{eq:mixing}. The
four names carrying an underscore-suffixed component follow the Annex A
pattern for a resolved vector quantity; the vane names have no Annex A
counterpart, this control effector being specific to the configuration,
and are formed to the same pattern.

\subsection{Constants}

Values are transcribed by the assembler from the geometry contract, never
typed by hand.

\begin{table}[htbp]
\centering
\caption{Constant variable definitions (contract 1.8.0 values).}
\label{tab:constants}
\begin{tabular}{@{}llll@{}}
\toprule
\textbf{varID} & \textbf{Value} & \textbf{Units} & \textbf{Meaning} \\
\midrule
\varid{WingAreaM2} & 0.1883 & m$^2$ & gross planform reference area $S$ \\
\varid{WingSpanM} & 1.0629 & m & reference span $b$ \\
\varid{WingChordM} & 0.1771 & m & reference chord $\bar c$ \\
\varid{DiskDiameterM} & 0.2030 & m & rotor diameter $D$ \\
\varid{XmrpM} & $-0.2423$ & m & moment reference point, FRD \\
\varid{YmrpM} & 0 & m & \\
\varid{ZmrpM} & 0 & m & \\
\bottomrule
\end{tabular}
\end{table}

\noindent
Parasite drag is deliberately \emph{not} among the constants; it is a
table, for the reason developed in Section~\ref{sec:drag}.

The reference planform is rectangular and unswept ($b = 1.0629$~m,
$\bar c = 0.1771$~m), giving aspect ratio $A\!R = 6.0$---a number that
recurs in the post-stall anchors of Section~\ref{sec:aerotables}. The
rotor is a three-bladed fan of diameter $0.2030$~m within a duct of bore
diameter $0.2090$~m and chord $0.09135$~m, centred at
$x_{\mathrm{frd}} = -0.46787$~m.

\subsection{Derived variables}
\label{sec:derived}

These are MathML \texttt{calculation} elements inside the file, evaluated
before any table lookup.

\begin{align}
\bar q &= \tfrac{1}{2}\,\rho\,V_\infty^2
  &&\text{dynamic pressure} \label{eq:qbar}\\[1mm]
J &= \frac{V_\infty}{n D}
  &&\text{advance ratio} \label{eq:J}\\[1mm]
\hat u &= \bigl(\cos\alpha\cos\beta,\ \sin\beta,\
  \sin\alpha\cos\beta\bigr)
  &&\text{free-stream unit vector, FRD} \label{eq:uhat}\\[1mm]
\alpha_{\mathrm{disk}} &= \arccos\bigl(\cos\alpha\,\cos\beta\bigr)
  &&\text{disk incidence (validity monitor)} \label{eq:alphadisk}\\[1mm]
\phi_w &= \operatorname{atan2}\!\bigl(\sin\beta,\
  \sin\alpha\cos\beta\bigr)
  &&\text{crossflow azimuth (diagnostic)} \label{eq:phiw}\\[1mm]
\sigma &= \arcsin\bigl(\sin\alpha\,\cos\beta\bigr)
  &&\text{total inclination} \label{eq:sigma}\\[1mm]
\hat p &= \frac{p\,b}{2V_\infty}, \quad
\hat q = \frac{q\,\bar c}{2V_\infty}, \quad
\hat r = \frac{r\,b}{2V_\infty}
  &&\text{reduced rates} \label{eq:rates}\\[1mm]
T_c &= \frac{T}{\bar q\,S}
  &&\text{interaction index} \label{eq:Tc}
\end{align}

Three of these warrant comment.

\paragraph{Equations~\eqref{eq:alphadisk} and \eqref{eq:phiw} are
monitors, not table indices.} With the rotor axis along body $+x$, the
angle between the free stream and the axis follows from the first
component of Eq.~\eqref{eq:uhat}, and $\phi_w$ is the azimuth of the
crossflow about that axis, zero in the pitch plane. Neither indexes a
table, for the reason given in Section~\ref{sec:axialonly}: the
propulsor model is axial-inflow only. They are computed and exported
precisely so that a consumer can \emph{detect} when it has left the
model's validity envelope, which is a more useful service than an axis
the solver cannot honestly fill.

\paragraph{Equation~\eqref{eq:sigma}} is the collapse variable of the
post-stall study: beyond separation the loading depends on attitude
almost entirely through $\sigma$, with cross-$\beta$ spread within 13\%
from $\alpha \approx 6\dg$ through $75\dg$. It is exported as a
diagnostic and is the natural coordinate in which to inspect the
post-stall tables for consistency.

\paragraph{Equation~\eqref{eq:Tc} closes no algebraic loop.} The
propulsor wrench does not depend on $T_c$, so thrust is available from
the \varid{prop*} tables before the \varid{coupling*} tables are
evaluated. The evaluation order is fixed: derived scalars, propulsor
tables, $T_c$, airframe tables, interaction tables, buildup.

%======================================================================
\section{The vane command space}
\label{sec:vanes}

\subsection{Why not four independent angles}

The duct-exit cruciform has four all-moving vanes, and the contract
states each one's hinge axis independently. Tabulating four independent
angles would require $7^4 = 2401$ combinations before the advance-ratio
axis is considered---infeasible, and structurally wrong, because a flight
controller does not command vanes individually. It commands moments.

\subsection{Mode shapes from the geometry}

All four hinge axes in the contract point radially \emph{outward}:
$[0,0,1]$ (bottom, since $z$ is down), $[0,-1,0]$ (left), $[0,0,-1]$
(top), and $[0,1,0]$ (right). Applying an equal right-hand-rule rotation
about four outward radial axes tilts every plate with the same handedness
as seen from outside the duct, so all four tangential forces circulate
the same way about the $x$ axis: the net force cancels by symmetry and a
pure rolling couple survives. Differential deflection of one opposed pair
instead redirects the jet transversely, producing a net force and the
associated moment. Hence the mixing
\begin{equation}
\begin{pmatrix}
\delta_{\mathrm{b}} \\ \delta_{\mathrm{l}} \\
\delta_{\mathrm{t}} \\ \delta_{\mathrm{r}}
\end{pmatrix}
=
\begin{pmatrix}
\phantom{-}0 & \phantom{-}1 & 1 \\
-1 & \phantom{-}0 & 1 \\
\phantom{-}0 & -1 & 1 \\
\phantom{-}1 & \phantom{-}0 & 1
\end{pmatrix}
\begin{pmatrix}
\delta_P \\ \delta_Y \\ \delta_R
\end{pmatrix},
\label{eq:mixing}
\end{equation}
in which the third column is the common mode (roll), the first is the
$y$-pair differential (pitch, redirecting the jet in $z$), and the second
is the $z$-pair differential (yaw, redirecting it in $y$). The
\emph{numeric} signs in Eq.~\eqref{eq:mixing}---which member of each pair
receives the positive angle---are a convention, and the file pins them by
static shot rather than by argument, for the reason given in
Section~\ref{sec:frames}.

\subsection{Three properties the tables must respect}

\paragraph{Each mode carries a full wrench.} Rotor swirl means the two
vanes of a differential pair meet the jet at different local incidence,
so a pitch command produces measurable yaw and roll components. These
cross-terms are recorded rather than assumed zero---see the generated
Table~\ref{tab:propmodes}, where they are visible. Each mode table also
carries $\Delta C_Q$, the back-pressure on the rotor, because the
rotor--vane coupling is two-way: vane blockage and upwash enter the
rotor's boundary condition, and the rotor's loading shifts measurably
when the vanes deflect.

\paragraph{Deflection is not antisymmetric in $\delta$.} The vane sits in
a jet whose dynamic pressure and swirl are set by the rotor, not by the
free stream, so the loads at $+\delta$ and $-\delta$ are not mirror
images. The tables therefore carry both signs explicitly; no mirroring is
applied.

\paragraph{Superposition was measured, and it partly fails.} The obvious
buildup---tabulate three modes, add them---was tested rather than
assumed, and the test changed the model. The result is
Section~\ref{sec:superposition}: pitch and yaw superpose to better than
$0.1\%$, but any pair involving roll is wrong by $6$ to $33\%$. The
mechanism is simple once seen, and it is the reason the test was worth
running.

\subsection{Travel limits}

The contract states mechanical stops of $\pm 20\dg$ and a soft limit of
$\pm 15\dg$. The tables are generated over $\pm 15\dg$, and an allocator
consuming this model should saturate the \emph{sum} of the three mode
commands at each individual vane---i.e. each row of
Eq.~\eqref{eq:mixing}---against the soft limit. That limit exists
precisely because a vane can be driven to its mechanical stop but should
not be in normal operation.

%======================================================================
\section{Breakpoint sets}

\begin{table}[htbp]
\centering
\caption{Breakpoint sets. Angles in degrees.}
\label{tab:breakpoints}
\begin{tabular}{@{}llp{78mm}@{}}
\toprule
\textbf{bpID} & \textbf{Count} & \textbf{Values} \\
\midrule
\varid{alphaBp} & 25 & $-4$, $-2$, 0, 2, 4, 6, 8, 10, 12, 14, 16, 18, 20,
22, 24, 26, 30, 35, 40, 45, 50, 60, 70, 80, 90 \\
\varid{betaBp} & 13 & $0, \pm5, \pm10, \pm15, \pm20, \pm25, \pm30$
(generated at $\beta \geq 0$, mirrored) \\
\varid{aileronBp} & 7 & $-20$, $-10$, $-5$, 0, 5, 10, 20 \\
\varid{jBp} & 11 & 0, 0.1, 0.2, 0.3, 0.4, 0.5, 0.6, 0.7, 0.8, 0.9, 1.0 \\
\varid{vaneBp} & 7 & $-15$, $-10$, $-5$, 0, 5, 10, 15 \\
\varid{tcBp} & 6 & 0, 0.5, 1, 2, 4, 8 \\
\bottomrule
\end{tabular}
\end{table}

The $\alpha$ grid is deliberately non-uniform: $2\dg$ spacing through the
stall break, where the anchored section model produces $C_{L\max}$ at
$\alpha = 18\dg$ and the loads change rapidly, coarsening to $5$--$10\dg$
in the plate regime where the loading varies slowly with attitude. The
$\beta$ grid exploits the unpowered airframe's mirror symmetry,
\begin{equation}
C_Y, C_l, C_n \ \text{odd in } \beta; \qquad
C_X, C_Z, C_m \ \text{even in } \beta,
\label{eq:mirror}
\end{equation}
a property confirmed numerically in the source sweeps rather than merely
assumed. The $T_c$ grid is logarithmically spaced because the induced
velocity ratio that drives the interaction varies as $\sqrt{T_c}$ in
momentum theory, so equal ratios of $T_c$ correspond to roughly equal
increments of effect.

%======================================================================
\section{Table inventory and provenance}
\label{sec:inventory}

Every table is listed. Each group states which solver path generated it,
so that any number in the model can be traced to a run.

\subsection{\varid{aero*}: airframe, power-off}
\label{sec:aerotables}

Source: the Level-2 coupled solve---vortex lattice with source-panelled
body and duct, sectional viscous feedback, and the anchored post-stall
section model (Kirchhoff attenuation driven by a computed separation
point, Viterna--Corrigan deep stall~\cite{viterna} with
aspect-ratio-aware $C_{d\max}$, and Rayleigh free-streamline centre of
pressure).

\begin{longtable}{@{}lllp{50mm}@{}}
\caption{Airframe tables (21).}
\label{tab:aerotables} \\
\toprule
\textbf{varID} & \textbf{Axes} & \textbf{Quantity} & \textbf{Notes} \\
\midrule
\endfirsthead
\multicolumn{4}{@{}l}{\small\itshape Table~\ref{tab:aerotables},
continued} \\
\toprule
\textbf{varID} & \textbf{Axes} & \textbf{Quantity} & \textbf{Notes} \\
\midrule
\endhead
\bottomrule
\endfoot
\varid{aeroCX} & $\alpha\times\beta$ & $C_X$ & baseline, controls
neutral, power-off \\
\varid{aeroCY} & $\alpha\times\beta$ & $C_Y$ & \\
\varid{aeroCZ} & $\alpha\times\beta$ & $C_Z$ & \\
\varid{aeroCl} & $\alpha\times\beta$ & $C_l$ & \\
\varid{aeroCm} & $\alpha\times\beta$ & $C_m$ & \\
\varid{aeroCn} & $\alpha\times\beta$ & $C_n$ & \\
\addlinespace
\varid{aeroCZq} & $\alpha$ & $\partial C_Z/\partial \hat q$ & from
\texttt{CL\_q\_nd}, sign per Table~\ref{tab:signs} \\
\varid{aeroCmq} & $\alpha$ & $\partial C_m/\partial \hat q$ & from
\texttt{Cm\_q\_nd} \\
\varid{aeroClp} & $\alpha$ & $\partial C_l/\partial \hat p$ & from
\texttt{Croll\_p\_nd} \\
\varid{aeroCnp} & $\alpha$ & $\partial C_n/\partial \hat p$ & from
\texttt{Cn\_p\_nd} \\
\varid{aeroClr} & $\alpha$ & $\partial C_l/\partial \hat r$ & from
\texttt{Croll\_r\_nd} \\
\varid{aeroCnr} & $\alpha$ & $\partial C_n/\partial \hat r$ & from
\texttt{Cn\_r\_nd} \\
\varid{aeroCYp} & $\alpha$ & $\partial C_Y/\partial \hat p$ & converted
by the assembler from the per-rad/s field \\
\varid{aeroCYr} & $\alpha$ & $\partial C_Y/\partial \hat r$ & likewise \\
\addlinespace
\varid{aeroDCX} & $\alpha\times\delta_a$ & $\Delta C_X$ & aileron
increment from the neutral \\
\varid{aeroDCY} & $\alpha\times\delta_a$ & $\Delta C_Y$ & configuration at
matched $\alpha$. \\
\varid{aeroDCZ} & $\alpha\times\delta_a$ & $\Delta C_Z$ & Swept one-sided
and mirrored, the \\
\varid{aeroDCl} & $\alpha\times\delta_a$ & $\Delta C_l$ & symmetry verified
numerically \\
\varid{aeroDCm} & $\alpha\times\delta_a$ & $\Delta C_m$ &
(Section~\ref{sec:aileron}); \\
\varid{aeroDCn} & $\alpha\times\delta_a$ & $\Delta C_n$ & these attenuate
through stall \\
\addlinespace
\varid{aeroCD0} & $\alpha$ & $C_{D_0}$ & body and duct parasite
\emph{only}, never the wing (Section~\ref{sec:drag}) \\
\end{longtable}

All rate-derivative tables are generated at $\beta = 0$ and carry the
taper of Section~\ref{sec:ratetaper}. The eight comprise the full set the
solver produces in reduced-rate form; $\partial C_X/\partial\hat q$ is
omitted because the underlying induced-drag derivative reflects induced
drag only and is not a trustworthy axial-force rate term.

\paragraph{The overlap decision.} Two solver paths cover
$\alpha \in [6\dg, 16\dg]$: the attached-flow matrix of Part
I~\cite{partone} and the coupled post-stall solve of Part
II~\cite{parttwo}. Tabulating one range from each would introduce a seam
in the middle of the most flight-critical region. The model takes the
Part II coupled solve as its single source across the whole $\alpha$
range, because that solve reduces to the attached solution below
separation by construction; the Part I matrix is retained as a
cross-check in the verification data rather than as table content. There
is therefore exactly one solver path behind every airframe number, and no
blend to justify.

\paragraph{Post-stall values are cycle means.} Beyond separation the
coupled fixed point does not converge to a steady state; it enters a
limit cycle, and the exported value is the cycle mean, flagged
unconverged in the source JSON with its residual. Independent iteration
paths agree on these means to four or more digits, so they are
reproducible quantities---but a consumer flying a smooth interpolation
through this region is modelling a regime that physically buffets. The
tables therefore carry DAVE-ML uncertainty bounds derived from the cycle
RMS, which is the mechanism the standard provides for exactly this
situation.

\paragraph{Anchors worth stating.} The anchored model produces
$C_{L\max} = 1.76$ at $\alpha = 18\dg$, emergent from the computed
separation tables rather than fitted; $C_N(90\dg) = 1.52$ against the
Viterna anchor
\begin{equation}
C_{d\max} = 1.11 + 0.018\,A\!R = 1.218 \quad (A\!R = 6),
\end{equation}
and a centre of pressure that walks from the attached quarter chord to
$x_{\mathrm{cp}} = 0.52\bar c$, landing on Rayleigh's free-streamline
mid-chord---a metric the model was never fitted to. $C_{L\max}$ is an
\emph{upper bound}: bubble bursting is not modelled, so the usable
incidence limit may be lower than the tabulated stall.

\paragraph{Rate derivatives past stall: the quantity does not exist.}
\label{sec:ratetaper}
This began as a declared assumption---the derivatives were to be tapered
linearly to zero over $\alpha \in [20\dg, 40\dg]$, on the reasoning that
holding the last attached value or zeroing discontinuously were both
worse. Measurement replaced the assumption with something firmer, and
the taper was never right.

A rate derivative is a linearization, and whether it holds is testable
for the price of one extra pair of solves: difference the coupled solve
at two perturbation amplitudes and compare. Doing so for roll damping
gives a sharp verdict.

\begin{table}[htbp]
\centering
\caption{Roll damping differenced at two roll-rate amplitudes. Either
the two agree to a tenth of a percent, or they disagree by tens to
hundreds---nothing sits between.}
\label{tab:ratelinearity}
\begin{tabular}{@{}rrrr@{}}
\toprule
$\alpha$ [deg] & coarse & fine & spread \\
\midrule
0 & $-0.545$ & $-0.545$ & 0.0\% \\
18 & $-0.374$ & $-0.374$ & 0.0\% \\
20 & $-0.168$ & $+0.255$ & 252\% \\
24 & $-0.159$ & $+0.063$ & 140\% \\
30 & $-0.086$ & $-0.115$ & 34\% \\
35 & $-0.129$ & $-0.193$ & 49\% \\
40 & $-0.233$ & $-0.234$ & 0.6\% \\
90 & $-0.425$ & $-0.425$ & 0.0\% \\
\bottomrule
\end{tabular}
\end{table}

Below $18\dg$ and above $40\dg$ the derivative is well defined. Between
$20\dg$ and $35\dg$ it is not one at all: at $\alpha = 20\dg$ the coarse
step returns $-0.168$ and the fine step $+0.255$, so the \emph{sign} is
set by the perturbation amplitude alone. Sections are crossing into and
out of stall across the perturbation, making the response non-smooth at
the scale of the difference rather than merely nonlinear.

So a taper, a hold, and a directly measured number are equally
fabrications there, because the tabulated quantity does not exist. The
model carries the two well-defined ranges and \emph{omits} the six
breakpoints between them: \varid{alphaRateBp} has 19 entries rather than
25, a lookup interpolates across the gap, and the file header says why.
That the boundary is so clean---0.1\% on one side, tens of percent on
the other---argues the band belongs to the flow rather than to the
tolerance chosen to detect it.

An earlier pass reported a sign reversal near $\alpha = 24$--$26\dg$ as
autorotation, from a single amplitude. It was not a property of the
aircraft but of how hard it was pushed, and the correction is recorded
here rather than quietly dropped.

\subsection{\varid{prop*}: rotor, duct, and vanes}

Source: the two-way rotor--vane coupled solve---a block Gauss--Seidel
alternation of the rotating-frame rotor-plus-duct solve and the
static-frame vane solve, with an angular-momentum budget on the swirl.
This is the method validated in the SciTech manuscript~\cite{scitech}.
Thrust is recovered as $-D_i$ and shaft torque as $-M_x$ of the
rotating-frame solve, both then converted by Eq.~\eqref{eq:export}.

\begin{longtable}{@{}lllp{46mm}@{}}
\caption{Propulsor tables (9). The vane tables are \emph{per-vane}, not
per-mode: one vane's response to its own deflection, serving all four
positions by the cruciform's rotational symmetry. See
Section~\ref{sec:superposition} for the measurement that forced this
structure.}
\label{tab:proptables} \\
\toprule
\textbf{varID} & \textbf{Axes} & \textbf{Quantity} & \textbf{Notes} \\
\midrule
\endfirsthead
\multicolumn{4}{@{}l}{\small\itshape Table~\ref{tab:proptables},
continued} \\
\toprule
\textbf{varID} & \textbf{Axes} & \textbf{Quantity} & \textbf{Notes} \\
\midrule
\endhead
\bottomrule
\endfoot
\varid{propCT} & $J$ & $C_T$ & vanes neutral; includes duct and vane
axial force \\
\varid{propCQ} & $J$ & $C_Q$ & shaft torque \\
\addlinespace
\varid{propDCXvane} & $J\times\delta_i$ & $\Delta C_X$ & one vane's own
increment, normalised per Eq.~\eqref{eq:propnorm} \\
\varid{propDCYvane} & $J\times\delta_i$ & $\Delta C_Y$ & \\
\varid{propDCZvane} & $J\times\delta_i$ & $\Delta C_Z$ & \\
\varid{propDClvane} & $J\times\delta_i$ & $\Delta C_l$ & \\
\varid{propDCmvane} & $J\times\delta_i$ & $\Delta C_m$ & \\
\varid{propDCnvane} & $J\times\delta_i$ & $\Delta C_n$ & \\
\varid{propDCQvane} & $J\times\delta_i$ & $\Delta C_Q$ & two-way
back-pressure on the shaft \\
\end{longtable}

\subsection{Axial inflow only: a declared structural limit}
\label{sec:axialonly}

The propulsor tables carry no disk-incidence axis. This is not an
economy; it is a statement about what the method can represent.

The rotor--vane coupling is built on axisymmetry at every stage. The
slipstream is reconstructed as radial bands of axial and swirl velocity;
the vanes' influence on the rotor enters as an azimuthally averaged
induced field; the self-consistent wake pitch is set by a radial inflow
distribution $v_i(r)$. Each of these assumes the loading is independent
of azimuth. Under non-axial inflow it is not: a blade advancing into the
crossflow meets a different relative velocity than one retreating, and
the resulting once-per-revolution loading produces the very in-plane
force and hub moment that a disk-incidence table would be created to
carry.

Generating an $\alpha_{\mathrm{disk}}$ sweep from this solver would
therefore produce numbers that vary---the free-stream projection changes
with incidence---while omitting the physical mechanism the axis exists to
capture. That is worse than an absent axis: it is a plausible artifact.
The model declares axial inflow, exports $\alpha_{\mathrm{disk}}$ as a
validity monitor (Eq.~\ref{eq:alphadisk}), and records non-axial
propulsor modelling as identified future work.

For a tail-sitter this limitation is real and bounded: the model is exact
in hover and in axial climb, and it degrades through the transition
corridor, where the free stream meets the disk at increasing incidence.
The interaction increments of Section~\ref{sec:coupling} are subject to
the same bound.

\subsection{\varid{coupling*}: fan-on-airframe increments}
\label{sec:coupling}

\begin{longtable}{@{}lllp{46mm}@{}}
\caption{Interaction tables (6), generated by the vortex-cylinder
induction model of Solver/DiskInduction.h.}
\label{tab:cpltables} \\
\toprule
\textbf{varID} & \textbf{Axes} & \textbf{Quantity} & \textbf{Notes} \\
\midrule
\endhead
\bottomrule
\endfoot
\varid{couplingDCX} & $\alpha\times T_c$ & $\Delta C_X$ & $\beta = 0$;
valid $V_\infty \geq 10$~m/s \\
\varid{couplingDCY} & $\alpha\times T_c$ & $\Delta C_Y$ & \\
\varid{couplingDCZ} & $\alpha\times T_c$ & $\Delta C_Z$ & the separation
delay appears here \\
\varid{couplingDCl} & $\alpha\times T_c$ & $\Delta C_l$ & \\
\varid{couplingDCm} & $\alpha\times T_c$ & $\Delta C_m$ & \\
\varid{couplingDCn} & $\alpha\times T_c$ & $\Delta C_n$ & \\
\end{longtable}

The mechanism is specific to this configuration and worth stating
precisely, because it is frequently misidentified. The ducted fan sits
\emph{aft} of the wing: the duct leading edge is $0.21$ chords behind the
wing trailing edge. The fan therefore does not blow the wing. What it
does is induce a favourable pressure gradient \emph{upstream} of itself,
over the wing, delaying separation---and it does so over precisely the
span that sheds first (duct outer radius to semi-span ratio $0.21$,
against a worst separation station at $|2y/b| \approx 0.15$--$0.23$).

The existing slipstream model cannot compute this: it is a
momentum-theory wake and returns identically zero upstream of the disk by
construction. Applied here it would report exactly zero interaction---an
artifact indistinguishable from a physical null result, and the most
dangerous kind of wrong answer because it is confidently null.

\paragraph{The prerequisite already existed.} These tables were long
recorded as blocked on an upstream-induction model. That was stale
bookkeeping rather than a missing capability: a uniformly loaded actuator
disk is exactly equivalent to a semi-infinite cylindrical vortex sheet,
which has a computable field everywhere \emph{including} ahead of the
disk and a closed form on the axis to check against, and the toolkit has
carried that model---tested, and already used by the attachment
sweep---for some time. Generating these tables required no new physics,
only a driver.

One feature of the configuration matters. The fan is an \emph{annulus}
around the tail boom, not a disk, and a uniformly loaded annulus sheds
trailing vorticity at both edges---the outer at $+\gamma_t$, the inner at
$-\gamma_t$. It is therefore two superposed cylinders, with the hub
radius taken from the body's own radius law at the duct station.

\paragraph{One sweep serves every speed.} Momentum theory gives
\begin{equation}
4\,\frac{A}{S}\left(1 + \frac{v_i}{V_\infty}\right)
\frac{v_i}{V_\infty} = T_c ,
\label{eq:tcsimilarity}
\end{equation}
so $v_i/V_\infty$ is a function of $T_c$ and the area ratio alone. The
induced field scales with $V_\infty$ and the \emph{coefficient}
increments therefore depend on $T_c$ only, which is what makes a
single-speed sweep valid across the envelope. The declared
$V_\infty \geq 10$~m/s limit is about $T_c$ degenerating as $\bar q$
vanishes, not about the field. The generated induced velocities satisfy
$2\rho A (V_\infty + v_i) v_i = T$ to every printed digit---the cheapest
available check that the disk loading was posed correctly.

% GENERATED by make-tables.py -- do not edit.
\begin{table}[htbp]
\centering
\caption{Fan-on-airframe lift increment $\Delta C_Z$ (negative is more lift) against incidence and thrust coefficient. The last column expresses the $T_c = 2$ result as a fraction of the power-off $C_Z$. Daggered rows are limit-cycle means in BOTH the powered and power-off solves, so those increments are differences of two cycle means.}
\label{tab:coupling}
\begin{tabular}{@{}rrrrrrr@{}}
\toprule
$\alpha$ [deg] & $T_c=0.5$ & $T_c=1$ & $T_c=2$ & $T_c=4$ & $T_c=8$ & vs.\ $C_Z$ \\
\midrule
0 & -0.0047 & -0.0076 & -0.0121 & -0.0186 & -0.0283 & +3.8\% \\
4 & -0.0072 & -0.0118 & -0.0186 & -0.0287 & -0.0434 & +2.7\% \\
8 & -0.0098 & -0.0160 & -0.0254 & -0.0391 & -0.0591 & +2.4\% \\
12 & -0.0132 & -0.0216 & -0.0340 & -0.0523 & -0.0790 & +2.3\% \\
14 & -0.0151 & -0.0247 & -0.0390 & -0.0597 & -0.0898 & +2.4\% \\
16 $^{\dagger}$ & -0.0523 & -0.0842 & -0.1319 & -0.1588 & -0.1939 & +7.5\% \\
20 $^{\dagger}$ & -0.0706 & -0.1052 & -0.1570 & -0.2334 & -0.3479 & +8.8\% \\
26 $^{\dagger}$ & -0.0486 & -0.0803 & -0.1284 & -0.2193 & -0.3294 & +7.5\% \\
35 $^{\dagger}$ & -0.0401 & -0.0664 & -0.1067 & -0.1680 & -0.2762 & +6.9\% \\
50 $^{\dagger}$ & -0.0254 & -0.0424 & -0.0686 & -0.1093 & -0.1731 & +4.7\% \\
70 $^{\dagger}$ & -0.0096 & -0.0160 & -0.0264 & -0.0434 & -0.0727 & +1.8\% \\
90 $^{\dagger}$ & 0.0119 & 0.0211 & 0.0793 & 0.0826 & 0.0980 & -5.2\% \\
\bottomrule
\end{tabular}
\end{table}


\noindent
The result splits at the stall. Below $\alpha \approx 14\dg$, where both
the powered and power-off solves converge cleanly, the fan adds
$2$--$4\%$ to the lift and the increment grows smoothly with incidence
and thrust alike. That is the accelerated stream over the wing, and it is
a solid number. From $16\dg$ the increment jumps to $7$--$9\%$, peaking
at $\alpha = 20\dg$---immediately past the configuration's own
$C_{L\max}$ at $18\dg$, which is precisely where a delayed separation
ought to pay. It then decays through deep stall and \emph{reverses sign}
by $90\dg$: with the plate broadside the fan's axial induction is
perpendicular to the free stream, so it slightly reduces the effective
incidence rather than energising an attached boundary layer.

\paragraph{What the jump does and does not establish.} It is tempting to
read the step at $16\dg$ as the separation-delay signature itself, and
the defensible position is weaker. The step coincides exactly with the
onset of limit-cycle behaviour in \emph{both} solves---residuals rise
from $10^{-4}$ to $0.2$--$0.6$ there---so every increment past it is a
difference of two cycle means, and the size of the step cannot be
separated from the change in solution character.

What does survive is that the increment scales cleanly and monotonically
with $T_c$ at every post-stall attitude: at $\alpha = 20\dg$ the ratios
between successive thrust levels are $1.49$, $1.49$ and $1.49$. Cycle-mean
scatter does not do that. So there is a real, thrust-ordered fan effect in
the stalled regime; attributing it specifically to delayed separation is
consistent with the data without being established by it.

The direct measurement is available, and is the right next step. The
coupled solve already computes a separation point on every strip, so a
driver can report the powered and power-off separation locations and
compare them, instead of inferring the delay from a lift increment. That
answers the question the fan-induction study actually poses---how far the
fan delays separation---without passing through a difference of two limit
cycles.

%======================================================================
\section{The buildup}
\label{sec:buildup}

The file assembles the total wrench in MathML, so that a consumer
receives forces and moments rather than a set of coefficients it must
combine correctly. Airframe coefficients accumulate baseline, control,
interaction and rate contributions:
\begin{align}
C_X &= \varid{aeroCX}(\alpha,\beta) + \varid{aeroDCX}(\alpha,\delta_a)
     + \varid{couplingDCX}(\alpha,T_c)
     - \varid{aeroCD0}(\alpha)\,\cos\alpha\cos\beta \\
C_Y &= \varid{aeroCY}(\alpha,\beta) + \varid{aeroDCY}(\alpha,\delta_a)
     + \varid{couplingDCY}(\alpha,T_c)
     + \varid{aeroCYp}\,\hat p + \varid{aeroCYr}\,\hat r
     - \varid{aeroCD0}(\alpha)\,\sin\beta \\
C_Z &= \varid{aeroCZ}(\alpha,\beta) + \varid{aeroDCZ}(\alpha,\delta_a)
     + \varid{couplingDCZ}(\alpha,T_c) + \varid{aeroCZq}\,\hat q
     - \varid{aeroCD0}(\alpha)\,\sin\alpha\cos\beta \\
C_l &= \varid{aeroCl}(\alpha,\beta) + \varid{aeroDCl}(\alpha,\delta_a)
     + \varid{couplingDCl}(\alpha,T_c)
     + \varid{aeroClp}\,\hat p + \varid{aeroClr}\,\hat r \\
C_m &= \varid{aeroCm}(\alpha,\beta) + \varid{aeroDCm}(\alpha,\delta_a)
     + \varid{couplingDCm}(\alpha,T_c) + \varid{aeroCmq}\,\hat q \\
C_n &= \varid{aeroCn}(\alpha,\beta) + \varid{aeroDCn}(\alpha,\delta_a)
     + \varid{couplingDCn}(\alpha,T_c)
     + \varid{aeroCnp}\,\hat p + \varid{aeroCnr}\,\hat r
\end{align}

Propulsor coefficients accumulate the baseline and the vane
contribution:
\begin{align}
C_X^{\mathrm{prop}} &= \varid{propCT}(J) + \Delta C_X^{\mathrm{vane}},
  \label{eq:propsumfirst}\\
C_k^{\mathrm{prop}} &= \Delta C_k^{\mathrm{vane}},
  \qquad k \in \{Y, Z, l, m, n\}, \\
C_Q^{\mathrm{prop}} &= \varid{propCQ}(J) + \Delta C_Q^{\mathrm{vane}}.
  \label{eq:propsumlast}
\end{align}

The vane contribution is \emph{not} a sum over the three command modes.
That structure was tested and rejected on measurement
(Section~\ref{sec:superposition}). It is instead a sum over the four
physical vanes, each evaluated at its own total commanded angle from
Eq.~\eqref{eq:mixing}:
\begin{equation}
\Delta C_k^{\mathrm{vane}} = \sum_{i \in \{\mathrm{b,l,t,r}\}}
  \varid{propDC}k\varid{vane}(J,\ \delta_i), \qquad
\delta_i = \bigl[\mathbf{M}\,(\delta_P, \delta_Y, \delta_R)^{\!\top}\bigr]_i ,
\label{eq:pervane}
\end{equation}
with $\mathbf{M}$ the mixing matrix of Eq.~\eqref{eq:mixing}. A single
per-vane table---one vane's response to its own deflection, the other
three neutral---serves all four positions by the cruciform's rotational
symmetry, with the wrench components permuted accordingly. This replaces
21 mode tables with 7 per-vane tables and, unlike the mode sum, respects
the measured physics.

Dimensionalisation then uses each group's own reference quantities, and
the two groups sum in body axes:
\begin{align}
F_x &= \bar q S\, C_X + \rho n^2 D^4\, C_X^{\mathrm{prop}} \\
F_y &= \bar q S\, C_Y + \rho n^2 D^4\, C_Y^{\mathrm{prop}} \\
F_z &= \bar q S\, C_Z + \rho n^2 D^4\, C_Z^{\mathrm{prop}} \\
M_x &= \bar q S b\, C_l + \rho n^2 D^5\, C_l^{\mathrm{prop}} \\
M_y &= \bar q S \bar c\, C_m + \rho n^2 D^5\, C_m^{\mathrm{prop}} \\
M_z &= \bar q S b\, C_n + \rho n^2 D^5\, C_n^{\mathrm{prop}}
\end{align}

Both groups' moments are already referred to the same point---the
generator exports the propulsor wrench about the contract's moment
reference point, not about the disk centre---so no transfer term appears
in the sum. Thrust and shaft torque are exported alongside as
\begin{equation}
T = \rho n^2 D^4\, C_X^{\mathrm{prop}}, \qquad
Q = \rho n^2 D^5\, C_Q^{\mathrm{prop}},
\end{equation}
because $T$ feeds Eq.~\eqref{eq:Tc} and $Q$ is what a powertrain model
needs.

%======================================================================
\section{First generated results}

The propulsor generator has been run against the 1.8.0 geometry contract
at the contract's reference rotor speed. The tables below are reproduced
verbatim from the generated JSON by \texttt{make-tables.py}---they are
included as evidence that the pipeline closes end to end, not as a
validated performance prediction, which awaits the anchor data identified
in~\cite{scitech}.

% GENERATED by make-tables.py -- do not edit.
\begin{table}[htbp]
\centering
\caption{Generated propulsor baseline: the contract's ducted fan at $n = 203.5$~rev/s ($\Omega = 1279$~rad/s), $D = 0.2030$~m, 3 blades, $\rho = 1.225$~kg/m$^3$, vanes neutral. Forces and moments in FRD about the contract's moment reference point. $Q$ is the shaft torque, i.e. the jet's angular-momentum flux.}
\label{tab:propbaseline}
\begin{tabular}{@{}rrrrrrrr@{}}
\toprule
$J$ & $V_\infty$ & $T$ & $Q$ & $C_T$ & $C_Q$ & $M_x$ & conv. \\
 & [m/s] & [N] & [N\,m] & & & [N\,m] & \\
\midrule
0.00 & 0.00 & 15.121 & 0.1632 & 0.1755 & 0.00933 & -0.0113 & yes \\
0.10 & 4.13 & 13.181 & 0.1729 & 0.1530 & 0.00988 & -0.0132 & yes \\
0.20 & 8.26 & 11.398 & 0.1754 & 0.1323 & 0.01003 & -0.0165 & yes \\
0.30 & 12.39 & 10.058 & 0.1787 & 0.1167 & 0.01022 & -0.0192 & yes \\
0.40 & 16.53 & 8.687 & 0.1767 & 0.1008 & 0.01010 & -0.0208 & yes \\
0.50 & 20.66 & 7.103 & 0.1643 & 0.0824 & 0.00939 & -0.0203 & yes \\
0.60 & 24.79 & 5.223 & 0.1390 & 0.0606 & 0.00795 & -0.0190 & yes \\
0.70 & 28.92 & 3.153 & 0.1002 & 0.0366 & 0.00573 & -0.0135 & yes \\
0.80 & 33.05 & 0.874 & 0.0453 & 0.0101 & 0.00259 & -0.0052 & yes \\
0.90 & 37.18 & -1.201 & -0.0217 & -0.0139 & -0.00124 & 0.2156 & yes \\
1.00 & 41.31 & -1.834 & -0.0971 & -0.0213 & -0.00555 & 0.7446 & yes \\
\bottomrule
\end{tabular}
\end{table}


\noindent
The behaviour is qualitatively correct for a ducted fan: thrust falls
monotonically with advance ratio as the blades unload, and the
cruciform's lateral force and moment components vanish at neutral
deflection, as Eq.~\eqref{eq:mixing} with $\delta_P = \delta_Y =
\delta_R = 0$ requires. The swirl-recovery vanes return the large
majority of the shaft torque, leaving a small residual net rolling
moment---which is the design intent of a swirl-recovery cruciform, and
the reason the aircraft's rolling moment cannot be read off the rotor
alone.

\paragraph{The two highest advance ratios are outside the envelope, and
show it.} Thrust crosses zero near $J \approx 0.85$ and both $C_T$ and
$C_Q$ are negative at $J = 0.9$ and $1.0$: the rotor has entered the
windmilling regime, which Section~\ref{sec:nondim} declares outside the
model's validity. The residual rolling moment simultaneously jumps by
more than an order of magnitude, from $-0.019$ to $+0.745$~N\,m, because
the swirl the vanes were sized to recover has reversed while the vanes
have not. These rows are retained in the generated map rather than
deleted---they are the measured location of the envelope boundary, and
knowing where the model stops behaving is worth more than a table that
ends without explanation---but the assembler truncates the \varid{jBp}
axis below the thrust reversal, and a consumer must not interpolate into
that corner. Modelling windmilling properly requires the
inflow-angle parameterisation noted in Section~\ref{sec:nondim}, not an
extension of this grid.

% GENERATED by make-tables.py -- do not edit.
\begin{table}[htbp]
\centering
\caption{Generated vane-mode increments at the $+15^\circ$ soft limit, taken from the neutral baseline at matched advance ratio. Each mode carries a full wrench: the off-axis components are the rotor swirl's signature, not numerical noise. $\Delta Q$ is the back-pressure on the shaft, which exists because the rotor--vane coupling is two-way.}
\label{tab:propmodes}
\begin{tabular}{@{}llrrrrrrr@{}}
\toprule
Mode & $J$ & $\Delta F_x$ & $\Delta F_y$ & $\Delta F_z$ & $\Delta M_x$ & $\Delta M_y$ & $\Delta M_z$ & $\Delta Q$ \\
 & & [N] & [N] & [N] & [N\,m] & [N\,m] & [N\,m] & [N\,m] \\
\midrule
Pitch & 0.00 & -0.130 & 0.001 & -1.085 & 0.0179 & -0.3181 & -0.0019 & 0 \\
 & 0.10 & -0.142 & 0.001 & -1.220 & 0.0191 & -0.3580 & -0.0019 & 0 \\
 & 0.20 & -0.183 & 0 & -1.513 & 0.0195 & -0.4444 & -0.0015 & 0 \\
 & 0.30 & -0.248 & 0 & -1.910 & 0.0196 & -0.5613 & -0.0011 & 0 \\
 & 0.40 & -0.325 & 0 & -2.361 & 0.0191 & -0.6941 & -0.0007 & 0 \\
 & 0.50 & -0.411 & 0 & -2.848 & 0.0177 & -0.8377 & -0.0005 & 0 \\
 & 0.60 & -0.534 & 0 & -3.578 & 0.0147 & -1.0525 & 0 & -0.0001 \\
 & 0.70 & -0.636 & 0 & -4.180 & 0.0107 & -1.2300 & 0 & -0.0002 \\
 & 0.80 & -0.751 & 0 & -4.900 & 0.0050 & -1.4420 & -0.0002 & -0.0004 \\
 & 0.90 & -0.869 & 0 & -5.854 & -0.0181 & -1.7220 & -0.0018 & -0.0006 \\
 & 1.00 & -0.794 & -0.003 & -6.233 & -0.0820 & -1.8282 & 0.0067 & -0.0009 \\
\addlinespace
Roll & 0.00 & -0.571 & 0.009 & -0.001 & -0.1333 & -0.0005 & -0.0025 & -0.0023 \\
 & 0.10 & -0.447 & 0.003 & -0.002 & -0.1508 & -0.0007 & -0.0007 & -0.0014 \\
 & 0.20 & -0.277 & -0.006 & -0.002 & -0.1911 & -0.0006 & 0.0017 & -0.0012 \\
 & 0.30 & -0.364 & -0.003 & -0.003 & -0.2474 & -0.0008 & 0.0008 & -0.0004 \\
 & 0.40 & -0.493 & -0.002 & -0.003 & -0.3134 & -0.0008 & 0.0007 & 0.0008 \\
 & 0.50 & -0.637 & -0.002 & -0.001 & -0.3890 & -0.0003 & 0.0005 & 0.0023 \\
 & 0.60 & -1.066 & -0.003 & 0 & -0.5127 & 0 & 0.0008 & -0.0047 \\
 & 0.70 & -1.097 & -0.003 & 0 & -0.5927 & -0.0001 & 0.0009 & 0.0035 \\
 & 0.80 & -1.431 & -0.002 & 0 & -0.7054 & -0.0001 & 0.0005 & 0.0011 \\
 & 0.90 & -1.695 & 0.001 & 0.001 & -0.8664 & 0.0002 & -0.0003 & -0.0044 \\
 & 1.00 & -1.630 & 0.003 & 0.004 & -1.0511 & 0.0012 & -0.0010 & -0.0077 \\
\addlinespace
Yaw & 0.00 & -0.066 & 1.075 & 0.001 & 0.0178 & 0.0019 & -0.3152 & 0.0003 \\
 & 0.10 & -0.140 & 1.220 & 0.001 & 0.0193 & 0.0019 & -0.3578 & 0 \\
 & 0.20 & -0.183 & 1.513 & 0 & 0.0201 & 0.0016 & -0.4442 & 0 \\
 & 0.30 & -0.248 & 1.910 & 0 & 0.0200 & 0.0012 & -0.5612 & 0 \\
 & 0.40 & -0.325 & 2.361 & 0 & 0.0195 & 0.0008 & -0.6940 & 0 \\
 & 0.50 & -0.411 & 2.848 & 0 & 0.0179 & 0.0005 & -0.8375 & 0 \\
 & 0.60 & -0.534 & 3.577 & 0 & 0.0148 & 0 & -1.0524 & -0.0001 \\
 & 0.70 & -0.636 & 4.179 & 0 & 0.0109 & 0 & -1.2297 & -0.0003 \\
 & 0.80 & -0.751 & 4.899 & 0 & 0.0051 & 0.0002 & -1.4418 & -0.0004 \\
 & 0.90 & -0.869 & 5.854 & 0 & -0.0182 & 0.0018 & -1.7220 & -0.0006 \\
 & 1.00 & -0.793 & 6.233 & -0.003 & -0.0826 & -0.0068 & -1.8280 & -0.0009 \\
\bottomrule
\end{tabular}
\end{table}


\subsection{The aileron: a silent zero, and where the flap actually
belongs}
\label{sec:aileron}

The model's open items included a routine-sounding check: confirm that
aileron roll authority decays past separation before shipping
\varid{aeroDCl}, since a table retaining attached-flow roll power at
$40\dg$ would model a simulator recovering from departures it should
not. The check was run, and it returned something worse than a bad decay
curve.

The same $10\dg$ antisymmetric deflection was solved three ways at every
attitude: on the contract's native eight-row chordwise lattice, on a
single-row lattice, and through the Level-2 coupled path that actually
generates the model's force tables.

\begin{table}[htbp]
\centering
\caption{Aileron rolling-moment increment for a $10\dg$ antisymmetric
command. Both single-row columns are exactly zero at every attitude.}
\label{tab:aileron}
\begin{tabular}{@{}rrrr@{}}
\toprule
$\alpha$ [deg] & 8 rows, inviscid & 1 row, inviscid & 1 row, coupled \\
\midrule
0 & 0.00925 & 0 & 0 \\
6 & 0.00936 & 0 & 0 \\
12 & 0.00926 & 0 & 0 \\
18 & 0.00896 & 0 & 0 \\
30 & 0.00781 & 0 & 0 \\
60 & 0.00318 & 0 & 0 \\
\bottomrule
\end{tabular}
\end{table}

\paragraph{Two reasonable contracts that cannot both be satisfied.} The
lattice builder splits a strip at the hinge line only when it has at
least two chordwise rows to split across; with one row it returns the
undivided strip and the deflection is never applied. The viscous
coupling, in turn, requires exactly one Weissinger row per strip, since
its whole premise is one section state per strip. Each rule is
defensible on its own. Together they make a hinged control surface
unrepresentable in the only solve that produces the model's tables, and
the failure is silent: the solve converges, reports sensible forces, and
returns an aileron increment of exactly zero.

Had the check not been run, the model would have shipped six
control-derivative tables of zeros---an aircraft with no roll control
whatsoever---and neither the generator, the assembler, nor the
DTD validation would have objected. Only a \texttt{checkData} shot
comparing a deflected condition against a nonzero expectation would have
caught it, which is an argument for generating verification data from a
\emph{different} path than the tables themselves. The file now carries
exactly such a shot, \varid{aileronRollAuthorityIsReal}.

\paragraph{A correction to the first reading of this measurement.} The
coupled column above was originally obtained by differencing
\texttt{res.Base.Croll}, which \texttt{SolveViscousCoupled} never
populates (Section~\ref{sec:drag}'s companion defect, and the reason
Section~\ref{sec:verification} insists on independent paths). That
column was therefore differencing two zeros, and reported ``no control
effect'' for a reason unrelated to the hinge. The finding survives
because it rests on the \emph{single-row inviscid} column, which uses
the plain lattice solve and is genuinely zero. Coefficients now go
through the tested extractor of \varid{Solver/BodyAxes.h}.

\paragraph{The inviscid column is not a substitute.} It is a real control
effect, but its decay---$0.00936$ at $\alpha = 6\dg$ falling to
$0.00318$ at $60\dg$---is purely geometric, arising from the free-stream
projection and induced effects in a lattice that contains no stall
anywhere. It passes through the configuration's own stall near $18\dg$
without a break of any kind. A table built from it would hand a
simulator most of its attached-flow roll authority deep into stall:
precisely the failure the original check was meant to prevent, arrived at
from the opposite direction.

\paragraph{The fix, and why it is not a workaround.} It is solver-side,
and it consists of not fighting the lattice. A hinge genuinely cannot be
represented on one chordwise row; but to a strip method a deflected flap
is not a geometric object at all, it is a \emph{camber change}. Thin
airfoil theory says so quantitatively: a plain flap hinged at $x_h$
shifts the section's zero-lift angle by $-\tau\delta$, with
\begin{equation}
\tau = 1 - \frac{\theta_h}{\pi} + \frac{\sin\theta_h}{\pi}, \qquad
\theta_h = \arccos(1 - 2x_h).
\label{eq:tau}
\end{equation}
A deflected strip therefore carries \varid{FlapChordFraction} and
\varid{FlapDeflectionDeg}, and every section model is posed against
\varid{EffectiveAlpha0Deg()} rather than the camber line's own
$\alpha_0$. The lattice continues to supply the induced field; the
section model, which is the authority on $c_l$ in the coupled solve,
now knows the flap is down. This is the standard strip-theory treatment
of a flapped wing, not an expedient.

The change is additive---with no flap set the accessor returns
$\alpha_0$ exactly---so every existing consumer is bit-identical, which
\texttt{TestViscousCoupling}, \texttt{TestPostStallSection} and
\texttt{TestRotorVaneCoupling} confirm.

\paragraph{It was validated, not asserted.} The same driver was re-run
after the fix, and the comparison it affords is the point: the
\emph{resolved-hinge} eight-row lattice is an independent geometric
representation of the very same flap.

\begin{table}[htbp]
\centering
\caption{Aileron rolling-moment increment after the fix, both columns in
FRD. The coupled solve reproduces the independent geometric
representation where both are valid, and departs from it exactly where
only one contains the physics.}
\label{tab:aileronfixed}
\begin{tabular}{@{}rrrr@{}}
\toprule
$\alpha$ [deg] & 8 rows, inviscid & coupled & ratio \\
\midrule
0 & $-0.009254$ & $-0.009384$ & 1.014 \\
6 & $-0.009361$ & $-0.008322$ & 0.889 \\
12 & $-0.009261$ & $-0.006620$ & 0.715 \\
20 & $-0.008815$ & $-0.004184$ & 0.475 \\
\bottomrule
\end{tabular}
\end{table}

At zero incidence the two agree to $1.4\%$. The ratio then falls
monotonically as separation develops, and by $\alpha = 30\dg$ only about
$12\%$ of the attached-flow roll authority survives, against the
inviscid lattice's $84\%$. That decay is precisely the physics the
inviscid column structurally cannot contain, and it is the answer to the
question this section opened with.

\paragraph{What the generated sweep then showed.} With the fix in place
the full $\alpha \times \delta_a$ sweep runs, and its result is worth
more than the table entries themselves.

% GENERATED by make-tables.py -- do not edit.
\begin{table}[htbp]
\centering
\caption{Aileron increments at $\delta_a = +10^\circ$ from the generated sweep. Roll authority collapses through stall while adverse yaw grows, so past $\alpha \approx 30^\circ$ the surface produces mostly yaw -- and its rolling moment is small enough that its sign is not resolved. Daggered rows are limit-cycle means, not steady states.}
\label{tab:aileronsweep}
\begin{tabular}{@{}rrrrl@{}}
\toprule
$\alpha$ [deg] & $\Delta C_l$ & $\Delta C_n$ & retained & $|\Delta C_n / \Delta C_l|$ \\
\midrule
0 & -0.009558 & +0.000871 & 100\% & 0.09 \\
6 & -0.009358 & +0.001037 & 98\% & 0.11 \\
12 & -0.008814 & +0.001158 & 92\% & 0.13 \\
18 & -0.007491 & +0.001411 & 78\% & 0.19 $^{\dagger}$ \\
20 & -0.006908 & +0.001484 & 72\% & 0.21 $^{\dagger}$ \\
26 & -0.005330 & +0.001908 & 56\% & 0.36 $^{\dagger}$ \\
30 & -0.000707 & +0.003281 & 7\% & 4.64 $^{\dagger}$ \\
45 & +0.000776 & +0.004061 & 8\% & 5.24 $^{\dagger}$ \\
60 & -0.001974 & +0.004810 & 21\% & 2.44 $^{\dagger}$ \\
90 & +0.000357 & -0.000001 & 4\% & 0.00 $^{\dagger}$ \\
\bottomrule
\end{tabular}
\\[1mm]\footnotesize $^{\dagger}$ limit-cycle mean.
\end{table}


\noindent
Roll authority falls to $72\%$ of its attached value by $\alpha =
20\dg$ and to under $10\%$ by $30\dg$. Adverse yaw moves the other way:
$|\Delta C_n / \Delta C_l|$ climbs from $0.09$ at zero incidence to
above unity past stall. \textbf{Beyond about $30\dg$ the aileron is
predominantly a yaw effector}, which is the classic pre-departure
signature and precisely the behaviour a departure-susceptible simulator
must have. A table built from the inviscid lattice would have shown
roll authority decaying gently and adverse yaw staying small---an
aircraft that rolls obediently out of a stall.

Two honesties about the deep-stall rows. They are limit-cycle means, so
the daggered entries are not steady states. And once $|\Delta C_l|$ has
collapsed to a few thousandths its \emph{sign} is not resolved by this
method---the sweep shows it changing sign between $30\dg$ and $60\dg$,
which is consistent with genuine aileron reversal but equally
consistent with cycle-mean scatter at that amplitude. The model
therefore carries the values with their uncertainty rather than a claim
about reversal.

\paragraph{The moment the lift shift leaves behind.} Equation
\eqref{eq:tau} is thin-airfoil theory's \emph{lift} result, and taking it
alone is not neutral. The same Glauert coefficients give the flap's
quarter-chord moment,
\begin{equation}
\Delta c_{m,c/4} = -\tfrac{1}{2}\,\delta\,\sin\theta_h\,
                    (1 - \cos\theta_h),
\label{eq:flapcm}
\end{equation}
which vanishes at both limits for reasons worth stating: a flap hinged at
the leading edge is the whole section rotating, adding incidence but no
camber, and a flap of zero chord does nothing. In between it does not
vanish. For the contract's $12\%$ aileron at $10\dg$ it is
$\Delta c_m = -0.100$ against $\Delta c_l = 0.474$, so the flap's load
acts $0.21$ chords \emph{aft} of the quarter chord---a nose-down couple
that a lift-only flap model simply loses. It is added to the section
model's own $c_m$ rather than replacing it, since the section knows its
camber, the flap changes that camber, and thin-airfoil theory is linear.

\paragraph{One number worth carrying away.} Equation~\eqref{eq:tau} is
strongly concave, so the contract's $12\%$-chord aileron is worth
$\tau = 0.432$, not $0.12$. A linear-in-chord estimate---the intuitive
one---understates roll authority by a factor of three and a half while
looking entirely reasonable. \texttt{TestFlapSection} pins $\tau$ at its
two exact limits ($\tau = 1$ for a leading-edge hinge, the whole section
rotating; $\tau = 0$ for zero flap chord), at the textbook half-chord
and quarter-chord values, and pins that a deflection produces a rolling
moment at all---the assertion that would have caught the original defect
the day it appeared.

\subsection{Parasite drag, and the case against a constant}
\label{sec:parasiteresults}

% GENERATED by make-tables.py -- do not edit.
\begin{table}[htbp]
\centering
\caption{Parasite drag of the body and duct, referred to the wing reference area $S = 0.1883$~m$^2$ at $V_\infty = 25$~m/s. Upper block: the component buildup. Lower block: the total against incidence, Eq.~\eqref{eq:cd0}. The wing is absent by design---its profile drag is already inside the force tables.}
\label{tab:parasite}
\begin{tabular}{@{}lrrrrr@{}}
\toprule
Component & $S_{\mathrm{wet}}$ [m$^2$] & $Re$ & $C_f$ & $FF$ & $C_{D_0}$ \\
\midrule
body & 0.1422 & $8.41 \times 10^{5}$ & 0.00462 & 1.493 & 0.00521 \\
duct & 0.1246 & $1.56 \times 10^{5}$ & 0.00649 & 1.181 & 0.00507 \\
\midrule
\multicolumn{5}{@{}l}{friction total, incl. 3\% excrescence allowance} & 0.01058 \\
\midrule
$\alpha$ [deg] & \multicolumn{2}{r}{friction} & \multicolumn{2}{r}{crossflow} & $C_{D_0}$ \\
\midrule
0 & \multicolumn{2}{r}{0.01058} & \multicolumn{2}{r}{0.00000} & 0.01058 \\
4 & \multicolumn{2}{r}{0.01058} & \multicolumn{2}{r}{0.00006} & 0.01065 \\
8 & \multicolumn{2}{r}{0.01058} & \multicolumn{2}{r}{0.00050} & 0.01108 \\
12 & \multicolumn{2}{r}{0.01058} & \multicolumn{2}{r}{0.00166} & 0.01224 \\
16 & \multicolumn{2}{r}{0.01058} & \multicolumn{2}{r}{0.00387} & 0.01445 \\
20 & \multicolumn{2}{r}{0.01058} & \multicolumn{2}{r}{0.00739} & 0.01798 \\
30 & \multicolumn{2}{r}{0.01058} & \multicolumn{2}{r}{0.02310} & 0.03369 \\
45 & \multicolumn{2}{r}{0.01058} & \multicolumn{2}{r}{0.06535} & 0.07593 \\
60 & \multicolumn{2}{r}{0.01058} & \multicolumn{2}{r}{0.12005} & 0.13063 \\
90 & \multicolumn{2}{r}{0.01058} & \multicolumn{2}{r}{0.18483} & 0.19541 \\
\bottomrule
\end{tabular}
\end{table}


\noindent
Two results are worth drawing out. First, the body and the duct
contribute almost equally to the friction branch ($0.00521$ against
$0.00507$), which is not obvious in advance: the duct's wetted area is
smaller, but its chord is short enough that its Reynolds number is a
factor of five lower and its skin-friction coefficient correspondingly
higher. A buildup that treated the duct as a minor appendage would lose
half the friction drag.

Second, and decisively for the model's structure, the crossflow branch
reaches $0.185$ at $90\dg$---\textbf{17.5 times} the friction total. Had
parasite drag been carried as the single constant that a cruise-envelope
model would naturally use, the tabulation would have omitted $95\%$ of
the parasite drag at the attitude this model exists to describe. Against
the configuration's $C_N(90\dg) = 1.52$ the parasite branch is a $13\%$
addition: not dominant, but far outside the tolerance any other number
here is quoted to. This is the measurement that converted
\varid{aeroCD0} from a constant into a table.

\paragraph{The duct's own separated drag.} The crossflow branch above is
a slender-body result and does not describe an annular ring meeting the
flow edge-on, so the shroud originally sat outside the incidence-dependent
term altogether and \varid{aeroCD0} was a declared lower bound at high
incidence---the wrong direction to be wrong in for a vehicle that
transits there. The ring now carries a term of its own on its
side-projected area, $2\,c_{\mathrm{duct}}(D_o - D_i) = 0.0029$~m$^2$,
which raises $C_{D_0}(90\dg)$ from $0.195$ to $0.208$.

That is deliberately a bluff-body \emph{estimate} and not a computation:
it borrows the crossflow coefficient and the finite-length factor from
the body rather than claiming independent knowledge of an annulus. The
ring is small, so the term is modest---a $6\%$ rise---but the tabulation
no longer omits an entire component from the branch that dominates at
transition attitudes. What remains unmodelled is the interference
between the body wake and the wing, so the estimate is still a lower
bound, by less than before.

\subsection{Vane-mode superposition: measured, and partly false}
\label{sec:superposition}

The natural buildup for a three-mode control system is to tabulate each
mode and add the increments, and that is what this model originally
specified. It was tested by solving each simultaneous pair of commands
directly and comparing against the sum of the corresponding single-mode
increments, all taken from the neutral baseline at matched advance ratio.
Each mode was commanded at $7.5\dg$ so that the per-vane sum---which is
what the mechanism actually sees---stays within the $\pm15\dg$ soft
limit.

% GENERATED by make-tables.py -- do not edit.
\begin{table}[htbp]
\centering
\caption{Measured superposition error of the vane command modes: each simultaneous pair against the sum of its two single-mode increments, both taken from the neutral baseline at matched advance ratio. Magnitudes are $\|F\| + 10\,\|M\|$ so that moment error is not swamped by force error. Pitch and yaw superpose; anything involving roll does not.}
\label{tab:superposition}
\begin{tabular}{@{}lrrrrr@{}}
\toprule
Pair & $\delta$ [deg] & $J$ & combined & error & relative \\
\midrule
pitch{+}yaw & +7.5 & 0.00 & 3.2015 & 0.0027 & 0.09\% \\
 & +7.5 & 0.20 & 4.4592 & 0.0021 & 0.05\% \\
 & +7.5 & 0.40 & 6.9067 & 0.0026 & 0.04\% \\
\addlinespace
pitch{+}roll & +7.5 & 0.00 & 1.9192 & 0.6182 & 32.21\% \\
 & +7.5 & 0.20 & 2.7324 & 0.7138 & 26.12\% \\
 & +7.5 & 0.40 & 4.5751 & 0.7666 & 16.76\% \\
\addlinespace
yaw{+}roll & +7.5 & 0.00 & 1.9114 & 0.6253 & 32.72\% \\
 & +7.5 & 0.20 & 2.7161 & 0.7136 & 26.27\% \\
 & +7.5 & 0.40 & 4.5637 & 0.7672 & 16.81\% \\
\bottomrule
\end{tabular}
\end{table}


\noindent
The result splits cleanly in two. \textbf{Pitch and yaw superpose
essentially exactly}, to within $0.04$--$0.09\%$ across every advance
ratio and both signs. \textbf{Any pair involving roll does not},
with errors from $6\%$ to $33\%$, worst in hover and at positive
commands.

The mechanism follows from the mixing matrix, Eq.~\eqref{eq:mixing}.
Pitch acts on the $y$-pair and yaw on the $z$-pair: \emph{disjoint}
sets of vanes. Commanding both leaves every vane at the same angle it
would have held under one command alone, so the loads add because
nothing about any individual vane has changed. Roll is the common mode
and deflects \emph{all four}---including the two a pitch or yaw command
is already deflecting. Under roll-plus-pitch the starboard vane sits at
$\delta_R + \delta_P = 15\dg$, whereas the two single-mode cases each
placed it at $7.5\dg$. A vane's load is not linear in its own
deflection, so the increments cannot add.

Two conclusions follow, and the second is the useful one:

\begin{enumerate}
\item \textbf{The mode-sum buildup is invalid} and was removed from
Section~\ref{sec:buildup}. Retaining it would have given a control
allocator a $33\%$ error in hover---precisely the flight condition in
which a tail-sitter depends on vane authority most.
\item \textbf{The failure is per-vane nonlinearity, not vane-to-vane
interference.} This is what the pitch-plus-yaw result establishes: that
case has all four vanes deflected simultaneously and still superposes to
$0.05\%$, so the mutual influence between vanes is negligible at these
deflections. What breaks superposition is only that a vane's own
response curves with its own angle.
\end{enumerate}

The second conclusion prescribes the replacement structure. If the
wrench is a sum of independent per-vane responses, each evaluated at that
vane's \emph{total} commanded angle, then the correct tabulation is
per-vane rather than per-mode---Eq.~\eqref{eq:pervane}. That structure is
nonlinear in each vane's angle by construction, and it is also cheaper:
seven tables in place of twenty-one, since one vane's response serves all
four positions under the cruciform's rotational symmetry.

\paragraph{The replacement was then verified, not assumed.} A separate
sweep deflected each vane alone, and the per-vane summation model was
used to predict every command in the superposition study---the three
single modes and the three pairs---from those single-vane responses
alone. Table~\ref{tab:pervane} compares the two candidate structures on
the same measurements.

\begin{table}[htbp]
\centering
\caption{Worst-case prediction error of the two candidate buildups,
over the three single modes and three simultaneous pairs at
$\pm 7.5\dg$ and $J \in \{0, 0.2, 0.4\}$.}
\label{tab:pervane}
\begin{tabular}{@{}lrr@{}}
\toprule
Buildup & worst error & tables required \\
\midrule
Sum of the three command modes & 32.7\% & 21 \\
Sum of the four per-vane responses, Eq.~\eqref{eq:pervane} & 1.4\% & 7 \\
\bottomrule
\end{tabular}
\end{table}

The per-vane model predicts every case to within $1.4\%$, and that worst
case is roll in hover---the same corner where the mode sum is worst by
$33\%$. Pitch and yaw are predicted to $0.13\%$. The model that is
simultaneously more accurate by a factor of twenty and cheaper by a
factor of three is the one the file adopts, and adopting it cost one
additional sweep of the driver that already existed.

\subsection{Three conditions that do not converge, and what was measured
about them}
\label{sec:nonconv}

Of the 209 conditions in the generated map, 206 met the outer iteration's
convergence tolerance of $0.01$. The three exceptions are all
\emph{positive} roll-mode commands at $J = 0.6$, with residuals of
$0.027$ to $0.066$. Because a reader is entitled to know whether such
cells are salvageable, both standard cures were tried and both failed:

\begin{table}[htbp]
\centering
\caption{The non-converging corner. Neither a larger iteration budget nor
a damped feedback gain resolves it.}
\label{tab:nonconv}
\begin{tabular}{@{}rrrrrr@{}}
\toprule
$\delta_R$ & \multicolumn{2}{c}{residual} & \multicolumn{2}{c}{thrust
[N]} & spread \\
\cmidrule(lr){2-3}\cmidrule(lr){4-5}
[deg] & 24 passes & 48 passes & $\omega = 0.35$ & $\omega = 0.15$ & \\
\midrule
$-15$ & 0.0061 & 0.0061 & \multicolumn{3}{c}{\itshape converged in 3
passes} \\
$-10$ & 0.0094 & 0.0094 & \multicolumn{3}{c}{\itshape converged in 2
passes} \\
$-5$ & 0.0066 & 0.0066 & \multicolumn{3}{c}{\itshape converged in 2
passes} \\
\addlinespace
$+5$ & 0.0656 & 0.0656 & 4.8766 & 4.8764 & 0.004\% \\
$+10$ & 0.0459 & 0.0465 & 4.5183 & 4.6779 & 3.5\% \\
$+15$ & 0.0268 & 0.0278 & 4.1554 & 4.1580 & 0.06\% \\
\bottomrule
\end{tabular}
\end{table}

Doubling the iteration budget reproduces the residuals to three digits,
and two of them worsen marginally---the signature of a stationary limit
cycle in the block Gauss--Seidel alternation, not of slow convergence.
Damping the vane-feedback relaxation from $0.35$ to $0.15$, the
conventional remedy for an oscillating Gauss--Seidel iteration, likewise
fails. The failure is moreover \emph{sign-asymmetric}: negative roll
commands converge in two or three passes at the same advance ratio. That
asymmetry is physical rather than numerical---a common-mode deflection
either adds to or opposes the residual swirl the vanes sit in, and only
one sign drives the vane row toward the incidence at which its own solve
is multivalued.

The useful outcome is not a fix but a measurement. Two independent
iteration paths bound the exported value: thrust agrees to $0.004\%$ and
$0.06\%$ at $\delta_R = +5\dg$ and $+15\dg$, so those cells are
reproducible despite missing tolerance, while $\delta_R = +10\dg$ differs
by $3.5\%$ and is genuinely path-dependent. The assembler carries all
three as DAVE-ML uncertainty bounds---the mechanism the standard provides
for precisely this---rather than as ordinary entries, and a consumer
therefore sees a stated confidence rather than a false precision.

%======================================================================
\section{Verification}
\label{sec:verification}

A generated file that nothing re-reads is the one artifact in a codebase
that can rot silently. The model is therefore verified by its own
contents.

\subsection{Static shots}

The standard provides for this directly: a \texttt{checkData} element
carries one or more \texttt{staticShot} children, each an input vector
(\texttt{checkInputs}), an optional dump of internal signals
(\texttt{internalValues}), and an output vector with per-signal matching
tolerances (\texttt{checkOutputs})~\cite{davemlref}. The model uses all
three, and the \texttt{internalValues} dump is what makes a failing shot
diagnosable rather than merely red---an intermediate coefficient that
already disagrees localises the fault to a table, while agreement
through to the final summation localises it to the buildup.

Every shot is generated from a solve; none is hand-written.

\begin{enumerate}
\item \textbf{Moment transfer.} One condition about the contract
reference point and about a point displaced $+0.10$~m in
$x_{\mathrm{frd}}$; the $C_m$ difference must satisfy
Eq.~\eqref{eq:transfer}. Detects the frame-conversion failure described
in Section~\ref{sec:frames}.
\item \textbf{Aileron sign.} $\delta_a = +5\dg$ at $\alpha = 4\dg$; the
solved $C_l$ defines the sign, and \varid{aileronDeg}'s variable
definition cites this shot.
\item \textbf{Rate derivative.} $C_{l_p} = -0.452$, the value already
pinned by the solver's unit tests against the textbook figure for an
$A\!R = 6$ wing.
\item \textbf{Vane modes.} Pure pitch, pure yaw and pure roll commands,
each pinning the dominant components and recording the off-axis
residuals at their solved values---these are physics, not zeros, and
pinning them at zero would be a falsification.
\item \textbf{Propulsor reference.} The SciTech reference case, whose
thrust already serves as that manuscript's own consistency anchor.
\item \textbf{Part I cross-checks.} Attached-range $C_L$ and $C_m$ at
several attitudes from the independent Part I matrix, at a tolerance
that accommodates the viscous-versus-inviscid difference. This is the
recorded justification for the overlap decision of
Section~\ref{sec:aerotables}.
\item \textbf{Summed wrench.} One fully populated condition---attitude,
rates, aileron, vane commands, powered---evaluated through the buildup of
Section~\ref{sec:buildup} to dimensional forces and moments, pinning
normalisation, summation and reference consistency end to end.
\end{enumerate}

\subsection{The evaluator}

Verification runs through a small in-repo interpreter: it parses the
file, performs multilinear interpolation on the gridded tables, evaluates
the MathML buildup, and compares each static shot against its stated
value. The alternative---depending on an external DAVE-ML
implementation---was rejected on the same grounds that keep the rest of
this codebase free of heavyweight dependencies, and because the MathML
subset the file uses is under the author's control precisely so that a
compact evaluator suffices. The verifier runs in continuous integration,
so the file cannot drift from the generators that produced it.

%======================================================================
\section{Generation pipeline}
\label{sec:pipeline}

\begin{lstlisting}
app/AeroMapExport.cpp        -> models/data/aero-map.json
app/PropulsionMapExport.cpp  -> models/data/propulsion-map.json
app/InductionMapExport.cpp   -> models/data/coupling-map.json   [blocked]
models/build-daveml.py       -> models/AetherionFlightModel.dml
models/verify-daveml.py      -> ctest suite
models/report/make-tables.py -> models/report/tables/*.tex
\end{lstlisting}

\noindent
Each driver exports JSON independently and the assembler consumes the
cached JSONs, so regenerating one map does not rerun the others. Rows are
flushed as they complete, so an interrupted long sweep retains every
finished condition. The JSON files are committed alongside the XML: the
model and the data that produced it version together. This report's data
tables are generated from the same JSON by the same rule---no number in
this document was typed by hand.

%======================================================================
\section{Declared limits}

Collected here so that a consumer need not reconstruct them from the body
of the report.

\begin{enumerate}
\item \textbf{Incompressible.} No Mach dependence exists in the
generating methods. Validity $M < 0.3$. The rotor tip Mach number at the
reference speed is $0.38$, so the blades themselves sit at the edge of
this assumption.
\item \textbf{Single Reynolds number.} All tables were generated at
$V_\infty = 25$~m/s, $\mathrm{Re}_n \approx 3\times 10^5$.
\item \textbf{Axial inflow on the propulsor}
(Section~\ref{sec:axialonly}): exact in hover and axial climb, degrading
through transition. No in-plane force or hub moment from disk incidence.
\item \textbf{Powered operation only.} The $\rho n^2 D^4$ group of
Eq.~\eqref{eq:propnorm} excludes $n \to 0$; windmilling and
low-rotor-speed descent are outside the envelope.
\item \textbf{Interaction increments past stall are differences of two
limit-cycle means} (Section~\ref{sec:coupling}). Below $14\dg$ both
solves converge and the increment is solid; above $16\dg$ neither does,
so the post-stall values carry a real thrust-ordered effect whose
attribution to delayed separation the data supports without
establishing. The interaction model is also uniform-loading, swirl-free,
and computed at $\beta = 0$ only.
\item \textbf{Aileron increments carry the flap as a section camber
change} (Section~\ref{sec:aileron}), not as panel geometry. The lift
shift is thin-airfoil theory's; the accompanying section pitching-moment
increment is not modelled, and neither is gap leakage nor the viscous
decay of effectiveness at large deflection---all of which reduce real
control power, so the tabulated authority is an upper bound.
\item \textbf{Parasite drag is a conceptual-design estimate}
(Sections~\ref{sec:drag} and~\ref{sec:parasiteresults}), carrying the
conventional $\pm10$--$20\%$ accuracy of a component buildup. The duct's
separated drag at incidence is now included as a bluff-body estimate on
the ring's side-projected area, raising $C_{D_0}(90\dg)$ from $0.195$ to
$0.208$; what remains unmodelled is the interference between the body
wake and the wing, so \varid{aeroCD0} is still a lower bound at high
incidence, by less than before.
\item \textbf{Rate derivatives are absent across
$\alpha \in [20\dg, 35\dg]$} (Section~\ref{sec:ratetaper}), because
measurement at two perturbation amplitudes shows no linear coefficient
exists there. The breakpoint axis omits those attitudes rather than
tapering, holding or fabricating a value.
\item \textbf{Post-stall values are cycle means} with uncertainty bounds,
not steady states.
\item \textbf{$C_{L\max}$ is an upper bound}: bubble bursting is not
modelled.
\item \textbf{No hysteresis.} The tables are the ascending-$\alpha$
branch. Static hysteresis is a known feature of this regime and a static
gridded table cannot express it; representation is deferred to a
dedicated study.
\item \textbf{Vane mode superposition} is spot-checked, not tabulated
pairwise.
\end{enumerate}

%======================================================================
\section{Open items}
\label{sec:open}

\begin{itemize}
\item Confirm the Annex A spellings in Table~\ref{tab:inputs} against
the published standard text rather than against the reference
documentation's examples, and settle the vane names, which have no
Annex A counterpart.
\item Validate the assembled file against the \texttt{DAVEfunc} DTD as a
build step, so that structural conformance is checked by a tool rather
than by reading.
\item Generate the airframe map (\varid{aero*}) and assemble the file.
\item Model the interference between the body wake and the wing, the
last contribution keeping \varid{aeroCD0} a lower bound at high
incidence.
\item Report the powered and power-off SEPARATION POINTS directly, which the coupled solve already computes per strip, instead of inferring the fan's separation delay from a lift increment taken across two limit cycles.
\item Extend the flap model with gap leakage and the viscous decay of
effectiveness at large deflection. The section pitching-moment
increment is now carried (Section~\ref{sec:aileron}); these two remain,
and both reduce real control power, so tabulated authority stays an
upper bound.
\item Measure and record vane-mode superposition error at soft-limit
corners.
\item Determine the sensitivity of the propulsor tables to the vane
azimuth assumption.
\item Confirm the rotor rotation sense against the aircraft's actual
installation; the generated map states the sense it used, and the
contract states the axis but not the handedness of the reference speed
about it.
\item Acquire the validation anchors identified in~\cite{scitech}
(isolated-propeller thrust and torque data; ducted-fan-with-vanes
experiment or RANS comparison) before any table is used for
certification-adjacent work.
\end{itemize}

%======================================================================
\begin{thebibliography}{99}

\bibitem{s119}
ANSI/AIAA S-119-2011 (reaffirmed 2016),
\emph{Flight Dynamics Model Exchange Standard},
American Institute of Aeronautics and Astronautics, Reston, VA, approved
22 March 2011. doi:10.2514/4.867965. Annex A, \emph{Standard Variable
Names}; Annex B, \emph{Dynamic Aerospace Vehicle Exchange Markup
Language (DAVE-ML) Reference}.

\bibitem{davemlref}
AIAA Modeling and Simulation Technical Committee, Simulation Standards
Subcommittee,
``Dynamic Aerospace Vehicle Exchange Markup Language (DAVE-ML)
Reference,'' Version 2.0.1, 31 March 2011.
\url{https://daveml.org/DTDs/prod/Ref/}

\bibitem{jackson2002}
Jackson, E. B., and Hildreth, B. L.,
``Flight Dynamic Model Exchange using XML,''
AIAA Modeling and Simulation Technologies Conference and Exhibit,
Monterey, CA, AIAA Paper 2002-4482, August 2002.
NASA NTRS 20030015249.

\bibitem{jackson2009}
Jackson, E. B., and Hildreth, B. L.,
``Benefits to the Simulation Training Community of a New ANSI Standard
for the Exchange of Aero Simulation Models,''
NASA Langley Research Center, NTRS 20090041825, 2009.

\bibitem{s119status}
``Status Report: ANSI/AIAA S-119-2011 Flight Dynamic Model Exchange
Standard,'' NASA Technical Reports Server, NTRS 20200004989, 2020.

\bibitem{r004}
ANSI/AIAA R-004-1992,
\emph{Recommended Practice for Atmospheric and Space Flight Vehicle
Coordinate Systems},
American Institute of Aeronautics and Astronautics, Washington, DC,
1992.

\bibitem{iso1151}
ISO 1151-1:1988,
\emph{Flight Dynamics---Concepts, Quantities and Symbols---Part 1:
Aircraft Motion Relative to the Air},
International Organization for Standardization, Geneva, 1988.

\bibitem{xml}
Bray, T., et al. (eds.),
\emph{Extensible Markup Language (XML) 1.1}, 2nd ed.,
W3C Recommendation, 2006.

\bibitem{mathml}
Carlisle, D., Ion, P., Miner, R., and Poppelier, N. (eds.),
\emph{Mathematical Markup Language (MathML) Version 2.0}, 2nd ed.,
W3C Recommendation, 2003.

\bibitem{partone}
Tuncer, O., \"U\c{c}ler, \c{C}., and G\"une\c{s}, A.,
``Aerodynamic Characterisation Across the Separation Boundary,
Part I: Attachment Lines, the Separation March, and the Attached-Flow
Envelope,'' manuscript, 2026.

\bibitem{parttwo}
Tuncer, O., \"U\c{c}ler, \c{C}., and G\"une\c{s}, A.,
``Aerodynamic Characterisation Across the Separation Boundary,
Part II: Anchored Post-Stall Sections and the Attitude Map to
$90\dg$,'' manuscript, 2026.

\bibitem{scitech}
Tuncer, O., \"U\c{c}ler, \c{C}., and G\"une\c{s}, A.,
``Ducted-Fan Thrust Vectoring with Downstream Control Vanes: A
Rotating-Frame Lattice Method with Two-Way Rotor--Vane Coupling,''
AIAA SciTech Forum, manuscript, 2026.

\bibitem{raymer2018}
Raymer, D. P.,
\emph{Aircraft Design: A Conceptual Approach}, 6th ed.,
AIAA Education Series, American Institute of Aeronautics and
Astronautics, Reston, VA, 2018.

\bibitem{hoerner1965fluiddynamicdrag}
Hoerner, S. F.,
\emph{Fluid-Dynamic Drag},
published by the author, Bricktown, NJ, 1965.

\bibitem{allenPerkins1951}
Allen, H. J., and Perkins, E. W.,
``A Study of Effects of Viscosity on Flow over Slender Inclined Bodies
of Revolution,'' NACA Report 1048, 1951.

\bibitem{jorgensen1977}
Jorgensen, L. H.,
``Prediction of Static Aerodynamic Characteristics for Slender Bodies
Alone and with Lifting Surfaces to Very High Angles of Attack,''
NASA TR R-474, 1977.

\bibitem{viterna}
Viterna, L. A., and Corrigan, R. D.,
``Fixed Pitch Rotor Performance of Large Horizontal Axis Wind
Turbines,'' NASA Lewis Research Center, Cleveland, OH, 1982.

\end{thebibliography}

\end{document}
